 % 本文件由 scripts/assemble.py 自动装配，勿手改（改 parts/ 后重跑）
% 第10章 弦凝聚——光与费米子的统一

% 第10章 STRING CONDENSATION—AN UNIFICATION OF LIGHT AND FERMIONS

\chapter{弦凝聚——光与费米子的统一}

\begin{keypoints}
\item \textbf{光与费米子是什么？} 光是任意大小的、由凝聚弦构成的网的涨落。费米子是凝聚弦的端点。
\item \textbf{光与费米子从何而来？} 光与费米子来自充满空间的凝聚弦网的集体运动。
\item \textbf{光与费米子为何存在？} 光与费米子之所以存在，是因为我们的真空恰好具有弦网凝聚（string-net condensation）。
\end{keypoints}

在第 9 章中，我们用投影构造在二维自旋系统中构造了许多量子有序态，并引入投影对称群（PSG）来表征和分类这些量子有序态。这些量子有序态不仅包含规范玻色子，而且还包含作为其低能集体激发的费米子。规范玻色子和费米子的演生是一个非常惊人的结果，因为底层的格子模型是一个纯粹的玻色模型。

许多年来，费米子和规范玻色子一直被视为基本的、不可触碰的。当有人提出费米子和规范理论可以从二维自旋液体中演生出来时（Baskaran et al., 1987; Kivelson et al., 1987; Baskaran and Anderson, 1988），这一提议甚至没有引起惊讶，因为它遭遇到的是怀疑与不信。这些怀疑是有根据的。即使到了现在，我们仍不清楚二维自旋液体是否存在。然而，这并不意味着最初的提议是错误的。把这种计算应用到 $SU(N)$ 自旋模型上时（Affleck and Marston, 1988），它的确在三维中给出良好定义的规范玻色子和费米子（见 10.7 节）（Wen, 2002a）。因此，费米子和规范玻色子并非那样基本而不可触碰。它们可以从卑微的玻色模型中的某些量子序演生出来。

然而，导致规范玻色子和费米子演生的投影构造是非常形式化的。似乎一切都依赖于一个非常不合情理的数学技巧——把一个自旋拆成两个非物理的费米子。人们很难对这样一种方法抱有任何信心。由这种方法得到的演生费米子和规范玻色子这一惊人结果，看起来实在难以置信。

本章可以被看作一个建立信心的练习。我们将讨论几个精确可解的自旋模型，它们同样可以用投影构造来求解。我们将证明，对那些模型，投影构造给出的结果与精确结果一致。特别地，正如投影构造所提示的，这些精确可解模型具有演生的费米子和一个规范场。

然而，更重要的是，精确可解模型把费米子和规范场的演生与一种现象——闭合弦网的凝聚——联系在了一起。我们发现，由投影构造得到的量子有序态实际上就是弦网凝聚态。表征由投影构造得到的态中的量子序的 PSG，实际上表征的是不同的弦网凝聚。投影构造与弦网凝聚只是描述同一类量子序的两种不同方式。弦网凝聚为形式化的投影构造提供了一个物理基础。

试图联系投影构造与弦网凝聚，只不过是本章诸讨论的一个形式上的动机。既然费米子和规范玻色子可以从弦网凝聚中演生，我们想为本章的讨论赋予一个更物理的动机。我们想问：弦网凝聚与我们宇宙中光和费米子的存在有没有关系？特别是，我们想就光与费米子提出下面这些问题。光与费米子是什么？光与费米子从何而来？光与费米子为何存在？

目前，对上述基本问题的标准回答似乎是“光是由规范场描写的粒子”以及“费米子是由反对易场描写的粒子”。在这里，我们想论证，对上述问题还存在另一种可能的（而且我相信是更物理的）回答，即：我们的真空充满了任意大小的、由弦状物体构成的网，而这些弦网形成一个量子凝聚态。按照弦网理论，光（以及其他规范玻色子）是凝聚弦网的集体振动，而费米子是弦的端点。我们看到，弦网凝聚为光和费米子提供了一个统一的起源。\footnote{这里，“弦网凝聚”指的是由任意大小的弦状物体构成的网的凝聚。}换句话说，弦网凝聚统一了光与费米子。如果有人说“要有弦网”，那么我们就同时得到了光和费米子。

在为光与费米子的弦网理论提供证据之前，我们想先澄清所谓“光存在”与“费米子存在”是什么意思。我们知道，物理中有一个天然的质量标度——普朗克质量。普朗克质量非常之大，以致我们看到的任何粒子，其质量至少比普朗克质量小 $10^{16}$ 倍。因此，与普朗克质量相比，所有观测到的粒子都可以当作无质量的来处理。当我们问某些粒子为什么存在时，我们的真正意思是：与普朗克质量相比，那些粒子为什么是无质量的（或近乎无质量的）。真正的问题是什么使某些类型的激发（例如光和费米子）成为无质量的（或近乎无质量的）。自然界究竟为什么想要拥有无质量的激发？是谁订购了它们？

其次，我们想澄清所谓“光与费米子的起源”是什么意思。我们知道，一切都必须来自某种东西。所以当我们问“光与费米子从何而来？”时，我们已经假定了存在比光和费米子更简单、更基本的东西。在 10.1 节中，我们将定义局域玻色模型（local bosonic model），它比具有规范场与费米子耦合的模型更简单。我们将把局域玻色模型看作更基本的。这样一种哲学将被称为局域性原理（locality principle）。我们将证明，如果一个局域玻色模型的基态中含有由弦状物体构成的网的凝聚，那么光和费米子就可以从这个模型中演生出来。

费米子的弦网理论解释了为什么我们的宇宙中费米子总是成偶数地出现——因为一条弦（或一个弦网）总有偶数个端点。关于规范玻色子和费米子的弦网理论还有一个实验预言，即所有费米子都必须携带某种规范电荷。乍看之下，这个预言似乎与已知的实验事实——中子不携带任何规范电荷——相矛盾。于是有人可能认为，关于规范玻色子和费米子的弦网理论已经被实验证伪了。这里我们想指出，如果我们假定我们的宇宙中还存在一种新的分立规范场，例如一个 $Z_2$ 规范场，那么关于规范玻色子和费米子的弦网理论仍然可以是正确的。在这种情形下，中子和中微子携带该分立规范场的一个非零电荷。因此，关于规范玻色子和费米子的弦网理论预言，我们的宇宙中存在分立的规范激发（例如规范通量线）。

按照量子序与弦网凝聚的图像，基本粒子（例如光子和电子）可能终究不是基本的。它们可能是普朗克标度以下某个局域玻色系统的集体激发。由于我们无法在接近普朗克标度处做实验，很难判定光子和电子究竟是不是基本粒子。在本章中，我们将通过二维和三维格子上的一些具体的局域玻色模型，证明光与费米子的弦网理论至少是自洽的。这里研究的局域玻色模型，只是长长的局域玻色模型清单中的几个例子，这些模型都包含演生的非禁闭费米子和规范玻色子（Foerster et al., 1980; Kalmeyer and Laughlin, 1987; Wen et al., 1989; Read and Sachdev, 1991; Wen, 1991a; Moessner and Sondhi, 2001; Balents et al., 2002; Ioffe et al., 2002; Motrunich and Senthil, 2002; Sachdev and Park, 2002; Kitaev, 2003; Motrunich, 2003; Wen, 2003c）。这些模型的基态都具有非平凡的拓扑序/量子序。

\begin{smallnote}
我们想指出，尽管有某些相似之处，上述关于规范玻色子和费米子的弦网理论不同于标准的超弦理论。在标准超弦理论中，闭合弦对应于引力子，开弦对应于规范玻色子。费米子来自世界面上的费米型场。超弦理论中所有基本粒子（包括规范玻色子）都对应于小弦的不同振动模式。而在弦网理论中，真空充满了大弦（或大弦构成的网，见图 1.2）。无质量规范玻色子对应于大闭合弦之网的涨落，费米子对应于开弦的端点。弦网理论中没有费米型场。

关于规范涨落的弦网理论与规范理论中的 Wegner--Wilson 圈密切相关（Wegner, 1971; Wilson, 1974; Kogut, 1979）。动力学规范理论与动力学 Wegner--Wilson 圈理论之间的关系由 Gliozzi et al.~(1979) 和 Polyakov (1979) 提出并研究。Savit (1980) 综述了规范理论与延展物体理论之间的各种对偶关系。特别地，人们发现某些统计格点规范模型对偶于某些统计膜模型（Banks et al., 1977）。这种对偶关系与量子模型中规范理论和弦网理论之间的关系直接相连。量子模型中的新特点是，弦的端点有时可以是费米子或任意子（Levin and Wen, 2003）。
\end{smallnote}

这里我们想强调，对于我们宇宙中实际的规范玻色子和费米子而言，弦网理论目前只是一个建议。尽管弦网凝聚能够产生并保护无质量的光子、胶子、夸克以及其他带电轻子，但我们目前还不知道弦网凝聚能否产生中微子——一种手征费米子。我们也不知道弦网凝聚能否产生弱相互作用的 $SU(2)$ 规范场——它与夸克和轻子手征地耦合。弦网凝聚在我们真空中的正确性，取决于上述问题的解决。

另一方面，如果我们只关心凝聚态物理的问题——如何用玻色子制造人工光和人工费米子——那么弦网理论和量子序确实给出了一个答案。要制造人工光和人工费米子，我们只需让某些弦凝聚。

\section{局域玻色模型}

在本章中，我们将只考虑局域玻色模型。我们认为局域玻色模型是基本的，因为它们是真正局域的。我们注意到，费米子模型一般是非局域的，因为不同格点上的费米子算符即使在格点相距很远时也不对易。与此相反，局域玻色模型是局域的，因为相距很远的玻色子算符彼此对易。由于局域玻色模型具有内在的局域性，我们相信自然的基本理论是一个局域玻色模型。为强调这一点，我们将给这一信念起一个名字：局域性原理。

让我们详细定义局域玻色模型。为了定义一个物理系统，我们需要明确（i）一个总希尔伯特空间，（ii）一组局域物理算符的定义，以及（iii）一个哈密顿量。基于这一理解，局域玻色模型定义为满足下列性质的模型：

\begin{enumerate}
\item[(i)] 总希尔伯特空间是有限维局域希尔伯特空间的直积。
\item[(ii)] 局域物理算符是作用在单个局域希尔伯特空间之内的算符，或是这些算符对邻近局域希尔伯特空间的有限乘积。我们也将称这些算符为局域玻色算符，因为相距很远时它们彼此对易。
\item[(iii)] 哈密顿量是局域物理算符之和。
\end{enumerate}

格子上的自旋 $1/2$ 系统是局域玻色模型的一个例子。其局域希尔伯特空间是二维的，包含 $|\uparrow\rangle$ 和 $|\downarrow\rangle$ 态。局域物理算符是 $\sigma_{\pmb{i}}^{a}$、$\sigma_{\pmb{i}}^{a}\sigma_{\pmb{i}+\pmb{x}}^{b}$ 等，其中 $\sigma^{a}$（$a=1,2,3$）是泡利矩阵。

自由的无自旋费米子系统（二维或更高维）不是局域玻色模型，尽管它拥有与自旋 $1/2$ 系统相同的总希尔伯特空间。这是因为不同格点上的费米子算符 $c_i$ 不对易，因而不是局域玻色算符。更重要的是，费米子跳跃算符 $c_i^\dagger c_j$ 在二维及更高维中不能写成局域玻色算符。（不过，由于 Jordan--Wigner 变换，一维费米子跳跃算符 $c_{i+1}^\dagger c_i$ 可以写成局域玻色算符。因此，如果把 $c_i$ 排除在局域物理算符的定义之外，一维费米系统可以是局域玻色模型。）格点规范理论不是局域玻色模型，因为它的总希尔伯特空间不能是局域希尔伯特空间的直积。

\section{由投影构造得到的一个精确可解模型}

在本节中，我们将构造一个精确可解的局域玻色模型。在下一节中，我们将证明该模型包含弦网凝聚和演生的费米子。

\subsection{精确可解模型的构造}

\begin{keypoints}
\item 精确可解模型可以通过寻找彼此对易的算符来构造。
\end{keypoints}

通常，投影构造并不给出精确结果。在本节中，我们将在二维正方格子上构造一个精确可解模型（Kitaev, 2003; Wen, 2003c）。我们的模型具有如下性质：投影构造给出精确的基态以及所有其他精确的激发态。构造的关键一步是找到一组彼此对易的算符。

让我们引入四个马约拉纳费米子（Majorana fermion）算符 $\lambda_{\pmb{i}}^{a}$（$a=x,\bar{x},y,\bar{y}$），满足 $\{\lambda_{a,\pmb{i}},\lambda_{b,\pmb{j}}\}=2\delta_{ab}\delta_{\pmb{i}\pmb{j}}$。我们发现，算符
\begin{equation*}
\hat U_{\pmb{i},\pmb{i}+\hat{\pmb{x}}}=\lambda_{\pmb{i}}^{x}\lambda_{\pmb{i}+\hat{\pmb{x}}}^{\bar{x}},\qquad
\hat U_{\pmb{i},\pmb{i}+\hat{\pmb{y}}}=\lambda_{\pmb{i}}^{y}\lambda_{\pmb{i}+\hat{\pmb{y}}}^{\bar{y}},\qquad
\hat U_{\pmb{i}\pmb{j}}=-\hat U_{\pmb{j}\pmb{i}},
\tag{10.2.1}
\end{equation*}
构成一组彼此对易的算符，即 $[\hat U_{\pmb{i}_1\pmb{i}_2},\hat U_{\pmb{j}_1\pmb{j}_2}]=0$。在得到一组对易算符之后，我们容易看出，作为诸 $\hat U_{\pmb{i}\pmb{j}}$ 的函数，相互作用费米子哈密顿量
\begin{equation*}
H=-\sum_{\pmb{i}}V_{\pmb{i}}\hat F_{\pmb{i}},\qquad
\hat F_{\pmb{i}}=\hat U_{\pmb{i},\pmb{i}_1}\hat U_{\pmb{i}_1,\pmb{i}_2}\hat U_{\pmb{i}_2,\pmb{i}_3}\hat U_{\pmb{i}_3,\pmb{i}},
\tag{10.2.2}
\end{equation*}
与所有 $\hat U_{\pmb{i}\pmb{j}}$ 对易。这里 $\pmb{i}_1=\pmb{i}+\hat{\pmb{x}}$，$\pmb{i}_2=\pmb{i}+\hat{\pmb{x}}+\hat{\pmb{y}}$，而 $\pmb{i}_3=\pmb{i}+\hat{\pmb{y}}$。我们将把 $\hat F_{\pmb{i}}$ 称为 $Z_2$ 通量算符。

为了得到式 (10.2.2) 中的哈密顿量 $H$ 所作用的希尔伯特空间，我们在每个格点上把 $\lambda^{x,\bar{x},y,\bar{y}}$ 组合成两个复费米子算符
\begin{equation*}
2\psi_{1,\pmb{i}}=\lambda_{\pmb{i}}^{x}+\mathrm{i}\lambda_{\pmb{i}}^{\bar{x}},\qquad
2\psi_{2,\pmb{i}}=\lambda_{\pmb{i}}^{y}+\mathrm{i}\lambda_{\pmb{i}}^{\bar{y}}
\tag{10.2.3}
\end{equation*}
可以验证，$\psi_{1,2}$ 满足标准的费米子反对易代数。这两个复费米子算符在每个格点上生成一个四维希尔伯特空间。总希尔伯特空间的维数为 $4^{N_{\text{site}}}$。诸 $\hat U_{\pmb{i}\pmb{j}}$ 与 $H$ 的对易性使我们能够精确地求解这个相互作用费米子系统。

为了得到 $H$ 的精确本征态和精确本征值，令 $|\{s_{ij}\}\rangle$ 为诸 $\hat U_{ij}$ 算符的具有本征值 $s_{ij}$ 的共同本征态。由于 $(\hat U_{ij})^{2}=-1$ 且 $\hat U_{ij}=-\hat U_{ji}$，我们发现 $s_{ij}$ 满足 $s_{ij}=\pm\mathrm{i}$ 且 $s_{ij}=-s_{ji}$。由于 $H$ 是诸 $\hat U_{ij}$ 的函数，我们发现 $|\{s_{ij}\}\rangle$ 也是式 (10.2.2) 的一个能量本征态，其能量为
\begin{equation*}
E=-\sum_{\pmb{i}}V_{\pmb{i}}F_{\pmb{i}},\qquad
F_{\pmb{i}}=s_{\pmb{i},\pmb{i}_1}s_{\pmb{i}_1,\pmb{i}_2}s_{\pmb{i}_2,\pmb{i}_3}s_{\pmb{i}_3,\pmb{i}}
\tag{10.2.4}
\end{equation*}
我们将把 $F_{\pmb{i}}$ 称为穿过方格 $\pmb{i}$ 的 $Z_2$ 通量。当 $V_{\pmb{i}}>0$ 时，基态由 $|\{s_{ij}\}\rangle$ 给出，其中 $s_{\pmb{i},\pmb{i}+\pmb{x}}=s_{\pmb{i},\pmb{i}+\pmb{y}}=\mathrm{i}$。对这样的态，我们有 $F_{\pmb{i}}=1$。为了得到激发态，我们翻转某些 $s_{ij}$ 的符号，使某些 $F_{\pmb{i}}=-1$。

为了考察诸 $|\{s_{ij}\}\rangle$ 是否代表 $H$ 的所有精确本征态，我们需要清点态的数目。设二维正方格子有 $N_{\text{site}}$ 个格点，且两个方向都取周期性边界条件。这时格子有 $2N_{\text{site}}$ 条链。由于 $s_{ij}$ 总共有 $2^{2N_{\text{site}}}$ 种不同的选择（每条链有两种选择），态 $|\{s_{ij}\}\rangle$ 遍尽了希尔伯特空间中全部 $4^{N_{\text{site}}}$ 个态。因此，诸 $\hat U_{\pmb{i}\pmb{j}}$ 的共同本征态不简并，上述方法使我们能够得到 $H$ 的所有本征态和本征值。我们已经精确求解了这个二维相互作用费米子系统！

我们注意到，哈密顿量 $H$ 只能把每个格点上的费米子数改变偶数。因此，$H$ 作用在一个每个格点具有偶数个费米子的子空间之内。我们将称这个子空间为物理希尔伯特空间。物理希尔伯特空间在每个格点上只有两个态。限制在物理空间之内时，$H$ 实际上描写一个自旋 $1/2$ 或硬核玻色子系统。为了得到相应的自旋 $1/2$ 哈密顿量，我们注意到
\begin{equation*}
\sigma_{\pmb{i}}^{x}=\mathrm{i}\lambda_{\pmb{i}}^{y}\lambda_{\pmb{i}}^{x},\qquad
\sigma_{\pmb{i}}^{y}=\mathrm{i}\lambda_{\pmb{i}}^{\bar{x}}\lambda_{\pmb{i}}^{y},\qquad
\sigma_{\pmb{i}}^{z}=\mathrm{i}\lambda_{\pmb{i}}^{x}\lambda_{\pmb{i}}^{\bar{x}}
\tag{10.2.5}
\end{equation*}
作用在物理希尔伯特空间之内，且满足泡利矩阵的代数。因此，我们可以把 $\sigma_{\pmb{i}}^{l}$ 认作自旋算符。利用在物理希尔伯特空间之内
\begin{equation*}
(-)^{\psi^{\dagger}_{1,\pmb{i}}\psi_{1,\pmb{i}}+\psi^{\dagger}_{2,\pmb{i}}\psi_{2,\pmb{i}}}
=\lambda_{\pmb{i}}^{x}\lambda_{\pmb{i}}^{y}\lambda_{\pmb{i}}^{\bar{x}}\lambda_{\pmb{i}}^{\bar{y}}=1
\end{equation*}
这一事实，我们可以证明，费米子哈密顿量 (10.2.2) 变为（见问题 10.2.2）
\begin{equation*}
H_{\text{spin}}=-\sum_{\pmb{i}}V_{\pmb{i}}\hat F_{\pmb{i}},\qquad
\hat F_{\pmb{i}}=\sigma_{\pmb{i}}^{x}\sigma_{\pmb{i}+\hat{\pmb{x}}}^{y}\sigma_{\pmb{i}+\hat{\pmb{x}}+\hat{\pmb{y}}}^{x}\sigma_{\pmb{i}+\hat{\pmb{y}}}^{y}
\tag{10.2.6}
\end{equation*}
在物理希尔伯特空间之内。

物理希尔伯特空间中的所有态（即自旋 $1/2$ 模型中的所有态）都可以从诸 $|\{s_{ij}\}\rangle$ 态投影到物理希尔伯特空间而得到，即 $\mathcal{P}|\{s_{ij}\}\rangle$。投影算符为
\begin{equation*}
\mathcal{P}=\prod_{\pmb{i}}\frac{1+(-)^{\psi^{\dagger}_{1\pmb{i}}\psi_{1\pmb{i}}+\psi^{\dagger}_{2\pmb{i}}\psi_{2\pmb{i}}}}{2}
\end{equation*}
由于 $[\mathcal{P},H]=0$，投影后的态 $\mathcal{P}|\{s_{ij}\}\rangle$（如果非零）仍然是 $H$（或 $H_{\text{spin}}$）的本征态，且具有相同的本征值。经过投影，相互作用费米子模型的精确解导致自旋 $1/2$ 模型的精确解。

我们注意到，对于 $x$ 和 $y$ 两个方向都有周期性边界条件的系统，所有链的乘积
\begin{equation*}
\prod_{\pmb{i}}\,(\mathrm{i}s_{\pmb{i},\pmb{i}+\pmb{x}})(\mathrm{i}s_{\pmb{i},\pmb{i}+\pmb{y}})=(-)^{\hat N_f},
\tag{10.2.7}
\end{equation*}
其中 $\hat N_f=\sum_{\pmb{i}}(\psi^{\dagger}_{1\pmb{i}}\psi_{1\pmb{i}}+\psi^{\dagger}_{2\pmb{i}}\psi_{2\pmb{i}})$ 是总费米子数算符。因此，$|\{s_{ij}\}\rangle$ 的投影非零，仅当
\begin{equation*}
\prod_{\pmb{i}}\,(\mathrm{i}s_{\pmb{i},\pmb{i}+\pmb{x}})(\mathrm{i}s_{\pmb{i},\pmb{i}+\pmb{y}})=1.
\tag{10.2.8}
\end{equation*}

物理态（每个格点具有偶数个费米子）在由
\begin{equation*}
\hat G=\prod_{\pmb{i}}G_{\pmb{i}}^{\;\psi^{\dagger}_{1,\pmb{i}}\psi_{1,\pmb{i}}+\psi^{\dagger}_{2,\pmb{i}}\psi_{2,\pmb{i}}}
\end{equation*}
生成的局域 $Z_2$ 变换下不变，其中 $G_{\pmb{i}}$ 是一个只取 $\pm1$ 两个值的任意函数。我们注意到，$Z_2$ 变换把 $\psi_{I\pmb{i}}$ 变为 $\tilde\psi_{I\pmb{i}}=G_{\pmb{i}}\psi_{I\pmb{i}}$，把 $s_{ij}$ 变为 $\tilde s_{ij}=G_{\pmb{i}}s_{ij}G_{\pmb{j}}$。我们发现，$|\{s_{ij}\}\rangle$ 与 $|\{\tilde s_{ij}\}\rangle$ 在投影后给出同一个物理态（如果它们的投影非零）。于是，$\{s_{ij}\}$ 是自旋 $1/2$ 系统 $H_{\text{spin}}$ 的精确本征态的一个多对一标签。向具有每格点偶数个费米子的物理希尔伯特空间的投影，使我们的理论成为一个 $Z_2$ 规范理论。

\subsection*{问题 10.2.1}

\begin{enumerate}
\item[(a)] 证明式 (10.2.1) 中的诸 $\hat U_{ij}$ 构成一组彼此对易的算符。
\item[(b)] 证明式 (10.2.3) 中的诸 $\psi$ 满足复费米子的标准反对易关系。
\end{enumerate}

\subsection*{问题 10.2.2}

\begin{enumerate}
\item[(a)] 证明 $2\psi_1^\dagger\psi_1-1=\mathrm{i}\lambda^{x}\lambda^{\bar{x}}$ 以及 $2\psi_2^\dagger\psi_2-1=\mathrm{i}\lambda^{y}\lambda^{\bar{y}}$。证明在一个格点上具有偶数个费米子的态是 $\lambda^{x}\lambda^{y}\lambda^{\bar{x}}\lambda^{\bar{y}}$ 的本征值为 $+1$ 的本征态。
\item[(b)] 证明式 (10.2.7)。
\item[(c)] 证明式 (10.2.5) 中的诸 $\sigma^{l}$ 在物理希尔伯特空间内满足泡利矩阵的代数。
\item[(d)] 证明费米子哈密顿量 (10.2.2) 在物理希尔伯特空间内变为自旋哈密顿量 (10.2.6)。
\end{enumerate}

\subsection*{问题 10.2.3}

我们已经看到，把费米子哈密顿量 $H$ 限制在每格点具有偶数个费米子的子空间之内，可以从费米子系统得到一个自旋 $1/2$ 系统。我们注意到，费米子哈密顿量 $H$ 也作用在每格点具有奇数个费米子的子空间之内。试把费米子哈密顿量 $H$ 限制在奇数费米子子空间之内，得到相应的自旋 $1/2$ 系统。

\subsection{精确本征态与拓扑简并的基态}

\begin{keypoints}
\item 基态简并受到 $Z_2$ 规范结构的保护，对任何局域微扰都是稳健的。
\end{keypoints}

上述 $Z_2$ 规范结构使我们能够清点由诸 $|\{s_{ij}\}\rangle$ 态投影得到的物理态的数目。我们再次假定两个方向都有周期性边界条件。注意到常数 $Z_2$ 规范变换 $G_{\pmb{i}}=-1$ 不改变 $s_{ij}$，我们看到，彼此规范等价的互不相同的 $s_{ij}$ 只有 $2^{N_{\text{site}}}/2$ 个。在 $4^{N_{\text{site}}}$ 个 $|\{s_{ij}\}\rangle$ 态中，有 $4^{N_{\text{site}}}/2$ 个满足 $\prod_{\pmb{i}}(\mathrm{i}s_{\pmb{i},\pmb{i}+\pmb{x}})(\mathrm{i}s_{\pmb{i},\pmb{i}+\pmb{y}})=1$（即具有偶数个费米子）。因此，诸 $|\{s_{ij}\}\rangle$ 态的投影至多给我们 $(4^{N_{\text{site}}}/2)\allowbreak/(2^{N_{\text{site}}}/2)\allowbreak=2^{N_{\text{site}}}$ 个物理态。另一方面，所有 $|\{s_{ij}\}\rangle$ 态的投影应当给我们全部 $2^{N_{\text{site}}}$ 个自旋态。因此，不同的 $s_{ij}$ 的 $Z_2$ 规范等价类必须导致独立的物理态，这样投影才能恢复全部 $2^{N_{\text{site}}}$ 个自旋态。$s_{ij}$ 的 $Z_2$ 规范等价类是周期格子上所有物理态的一对一标签。这样，我们能够得到周期格子上 $H_{\text{spin}}$ 的所有本征态和本征值。

我们注意到，$Z_2$ 通量 $F_{\pmb{i}}$ 在 $Z_2$ 规范变换下不变。因此，如果两个位形 $s_{ij}$ 与 $\tilde s_{ij}$ 分别具有不同的 $Z_2$ 通量 $F_{\pmb{i}}$ 与 $\tilde F_{\pmb{i}}$，则 $s_{ij}$ 与 $\tilde s_{ij}$ 必定属于两个不同的 $Z_2$ 规范等价类。两个态 $|s_{ij}\rangle$ 与 $|\tilde s_{ij}\rangle$ 在投影后将导致不同的物理态。这一结果提示，投影后的态（物理态）用 $Z_2$ 通量 $F_{\pmb{i}}$ 来标记更好。结果发现，$F_{\pmb{i}}$ 并不是物理态的一个完美的一对一标签。对于周期性边界条件下的偶 $\times$ 偶格子上的系统，每个 $F_{\pmb{i}}$ 标记 4 个态（即对每一个 $Z_2$ 通量位形 $F_{\pmb{i}}$，有且仅有 4 个物理态给出同一个 $F_{\pmb{i}}$）。具有相同 $F_{\pmb{i}}$ 的 4 个物理态全部具有相同的能量。

为了理解这一四重简并，让我们考虑由某个 $|\{s_{ij}\}\rangle$ 态的投影所给出的一个本征态。其他简并的本征态可以通过施行下列两个变换得到：
\begin{align*}
T_1&:(s_{\pmb{i},\pmb{i}+\pmb{x}},s_{\pmb{i},\pmb{i}+\pmb{y}})\to(s_{\pmb{i},\pmb{i}+\pmb{x}},(-)^{\delta_{iy}}s_{\pmb{i},\pmb{i}+\pmb{y}})\\
T_2&:(s_{\pmb{i},\pmb{i}+\pmb{x}},s_{\pmb{i},\pmb{i}+\pmb{y}})\to((-)^{\delta_{ix}}s_{\pmb{i},\pmb{i}+\pmb{x}},s_{\pmb{i},\pmb{i}+\pmb{y}})
\end{align*}
我们看到，$T_1$ 不改变 $s_{\pmb{i},\pmb{i}+\pmb{x}}$，只在 $i_y=0$ 时翻转 $s_{\pmb{i},\pmb{i}+\pmb{y}}$ 的符号。由构造可知，$T_1$ 与 $T_2$ 不改变 $Z_2$ 通量 $F_{\pmb{i}}$。如果我们把周期格子看作一个环面，那么 $T_1$ 与 $T_2$ 通过环面的两个洞各插入 $\pi$ 通量（见图 9.6）。

在偶 $\times$ 偶格子上，变换 $T_1$ 与 $T_2$ 不改变乘积 $\prod_{\pmb{i}}(\mathrm{i}s_{\pmb{i},\pmb{i}+\pmb{x}})(\mathrm{i}s_{\pmb{i},\pmb{i}+\pmb{y}})$。因此，三个变换 $T_1$、$T_2$ 与 $T_1T_2$ 生成另外三个简并态。我们看到，所有能量本征值都具有四重简并。特别地，自旋 $1/2$ 模型 $H_{\text{spin}}$ 在偶 $\times$ 偶周期格子上具有 4 个简并基态。

在偶 $\times$ 奇格子上，情况稍有不同。$T_2$ 生成的态具有奇数个费米子，不对应于任何物理自旋 $1/2$ 态。因此，我们只能用 $T_1$ 来生成另一个简并态。在偶 $\times$ 奇周期格子上只有两个简并基态（由 $T_1$ 生成）。在奇 $\times$ 奇格子上，也有由 $T_1T_2$ 生成的两个简并基态。

我们注意到，局域地看，$T_1$ 与 $T_2$ 变换与 $Z_2$ 规范变换不可区分。由于物理自旋算符在 $Z_2$ 规范变换下不变，它们在 $T_1$ 与 $T_2$ 变换下也不变。因此，即使我们在精确可解模型 (10.2.6) 上加入任意局域微扰，由 $T_1$ 与 $T_2$ 生成的简并基态仍将保持简并。基态的这一简并是一个稳健的拓扑性质，表明基态中存在非平凡的拓扑序。

\subsection*{问题 10.2.4}

考虑 $L_x\times L_y$ 周期格子上的式 (10.2.6)。我们假定 $V_{\pmb{i}}=V$，只在某一行方格上 $V_{\pmb{i}}=0$。这时可以把该系统看作定义在一个圆柱上，圆柱在 $x$ 方向具有两条圆形边缘。试求系统的基态简并度。这些简并基态可以看作边缘态。证明每个边缘格点有 $\sqrt{2}$ 个边缘态。这意味着边缘激发由一个马约拉纳费米子描写。

\subsection{基态的投影对称群表征}

\begin{keypoints}
\item 式 (10.2.6) 的 $V<0$ 与 $V>0$ 基态具有不同的量子序。
\end{keypoints}

在本节中，我们将假定自旋 $1/2$ 系统 (10.2.6) 中的 $V_{\pmb{i}}$ 是均匀的，即 $V_{\pmb{i}}=V$。我们已经通过把自旋算符写成费米子算符的乘积求解了式 (10.2.6)（见式 (10.2.5)），也就是说，我们用投影构造求解了自旋 $1/2$ 系统。十分有趣的是，对于这个特定的自旋 $1/2$ 系统 (10.2.6)，投影构造竟然给出精确的结果。由于精确基态波函数是由投影构造得到的（即通过投影自由费米子波函数 (9.2.11) 而得到），我们可以用第 9 章发展的 PSG 表征方法来表征基态中的量子序。

第 9 章的讨论限于自旋旋转不变的态。下面，我们将把平均场理论和 PSG 推广到自旋旋转非不变的态。如果我们把 $|\downarrow\rangle$ 态认作零玻色子态 $|0\rangle$，把 $|\uparrow\rangle$ 态认作单玻色子态 $|1\rangle$，那么自旋 $1/2$ 模型 $H_{\text{spin}}$ 也可以看作一个硬核玻色子模型。下面我们将用玻色子图像来描写我们的模型。

为了用投影构造构造量子有序（或纠缠）的多玻色子波函数，我们首先引入“平均场”费米子哈密顿量
\begin{equation*}
H_{\text{mean}}=\sum_{\langle ij\rangle}\left(\psi^{\dagger}_{I,\pmb{i}}u_{ij}^{IJ}\psi_{J,\pmb{j}}+\psi^{\dagger}_{I,\pmb{i}}w_{ij}^{IJ}\psi^{\dagger}_{J,\pmb{j}}+\text{h.c.}\right)
\tag{10.2.9}
\end{equation*}
其中 $I,J=1,2$。我们将用 $u_{ij}$ 和 $w_{ij}$ 表示元素分别为 $u_{ij}^{IJ}$ 和 $w_{ij}^{IJ}$ 的 $2\times2$ 复矩阵。令 $|\Psi^{(u_{ij},w_{ij})}_{\text{mean}}\rangle$ 为上述自由费米子哈密顿量的基态，则下列多体玻色子波函数可以得到：
\begin{equation*}
\Phi^{(u_{ij},w_{ij})}(\pmb{i}_1,\pmb{i}_2...)=\langle0|\prod_{n}b(\pmb{i}_n)|\Psi^{(u_{ij},w_{ij})}_{\text{mean}}\rangle
\tag{10.2.10}
\end{equation*}
其中
\begin{equation*}
b(\pmb{i})=\psi_{1,\pmb{i}}\psi_{2,\pmb{i}}
\tag{10.2.11}
\end{equation*}
我们注意到，物理玻色子波函数 $\Phi^{(u_{ij},w_{ij})}(\{\pmb{i}_n\})$ 在下列 $SU(2)$ 规范变换下不变：
\begin{equation*}
(\psi_{\pmb{i}},u_{ij},w_{ij})\to(G(\pmb{i})\psi_{\pmb{i}},\,G(\pmb{i})u_{ij}G^{\dagger}(\pmb{j}),\,G(\pmb{i})w_{ij}G^{\top}(\pmb{j}))
\end{equation*}
其中 $G(\pmb{i})\in SU(2)$。

玻色子波函数 $\Phi^{(u_{ij},w_{ij})}(\{\pmb{i}_n\})$ 中的量子序可以用 PSG 来表征。PSG 由使拟设 $(u_{ij},w_{ij})$ 不变的对称变换与 $SU(2)$ 规范变换的联合变换构成。

这里我们想指出，诸 $\hat U_{ij}$ 的共同本征态，即 $|\{s_{ij}\}\rangle$，是自由费米子系统
\begin{equation*}
\tilde H_{\text{mean}}=\sum_{\langle ij\rangle}\left(s_{ij}\hat U_{ij}+\text{h.c.}\right)
\tag{10.2.12}
\end{equation*}
的基态。因此，从投影构造的观点看，$\tilde H_{\text{mean}}$ 可以看作“平均场”哈密顿量 $H_{\text{mean}}$。事实上，$\tilde H_{\text{mean}}$ 是 $H_{\text{mean}}$ 的一个特例，其中
\begin{equation*}
-w_{i,i+\hat{\pmb{x}}}=u_{i,i+\hat{\pmb{x}}}=-\mathrm{i}s_{i,i+\hat{\pmb{x}}}(1+\tau^{3}),\qquad
-w_{i,i+\hat{\pmb{y}}}=u_{i,i+\hat{\pmb{y}}}=-\mathrm{i}s_{i,i+\hat{\pmb{y}}}(1-\tau^{3})
\end{equation*}
如果 $s_{ij}$ 与 $(u_{ij},w_{ij})$ 通过上式相联系，则态 $|\{s_{ij}\}\rangle$ 就等于平均场态 $|\Psi^{(u_{ij},w_{ij})}_{\text{mean}}\rangle$。物理自旋波函数由平均场态的投影 $\mathcal{P}|\{s_{ij}\}\rangle$ 得到，这与式 (10.2.10) 等价。

记住，在投影构造中，我们需要选择拟设 $u_{ij}$ 与 $w_{ij}$，使平均能量 $\langle\Psi^{(u_{ij},w_{ij})}_{\text{mean}}|\mathcal{P}H_{\text{spin}}\mathcal{P}|\Psi^{(u_{ij},w_{ij})}_{\text{mean}}\rangle$ 极小。对我们的自旋模型 $H_{\text{spin}}$，这样得到的平庸基态恰好就是精确基态！而且，如果我们选择一个不同的 $s_{ij}$，投影态 $\mathcal{P}|\Psi^{(u_{ij},w_{ij})}_{\text{mean}}\rangle$ 将对应于 $H_{\text{spin}}$ 的一个具有能量 (10.2.4) 的精确本征态。正是在这个意义上，投影构造提供了 $H_{\text{spin}}$ 的一个精确解。

当 $V>0$ 时，我们模型的基态由 $Z_2$ 通量位形 $F_{\pmb{i}}=1$ 给出。为了产生这样的通量，我们可以选择 $s_{\pmb{i},\pmb{i}+\pmb{x}}=s_{\pmb{i},\pmb{i}+\pmb{y}}=\mathrm{i}$。这时，式 (10.2.12) 变为带有
\begin{equation*}
-w_{i,i+\pmb{x}}=u_{i,i+\pmb{x}}=1+\tau^{3},\qquad
-w_{i,i+\pmb{y}}=u_{i,i+\pmb{y}}=1-\tau^{3}.
\end{equation*}
的式 (10.2.9)。为得到上述拟设的 IGG，我们注意到 $u_{ij}$ 在常数规范变换 $G(\pmb{i})=\mathrm{e}^{\mathrm{i}\phi\tau^{z}}$ 下不变。然而，$w_{ij}$ 仅当 $\phi=0,\pi$ 时才不变。于是，IGG $=Z_2$。我们发现低能有效理论是一个 $Z_2$ 规范理论。由于该拟设已经是平移不变的，它在带平凡规范变换 $G_x=G_y=\tau^{0}$ 的 $G_xT_x$ 与 $G_yT_y$ 下不变。上述拟设的 PSG 就是式 (9.4.10) 中的 Z2A PSG。因此，$V>0$ 的基态是一个 Z2A 态。注意，对 Z2A PSG，$G_xT_x$ 与 $G_yT_y$ 满足平移代数 $(G_xT_x)^{-1}(G_yT_y)^{-1}(G_xT_x)(G_yT_y)=1$。

当 $V<0$ 时，基态由位形 $F_{\pmb{i}}=-1$ 给出，它可以由 $s_{\pmb{i},\pmb{i}+\pmb{x}}=(-)^{i_y}s_{\pmb{i},\pmb{i}+\pmb{y}}=\mathrm{i}$ 产生（译注：按本页下方拟设中 $y$ 方向链上的 $(-)^{i_x}$ 因子以及 $F_{\pmb{i}}=-1$ 的要求，此处指数应为 $i_x$，即 $s_{\pmb{i},\pmb{i}+\pmb{y}}=(-)^{i_x}\mathrm{i}$；原文印作 $(-)^{i_y}$）。拟设现在具有形式
\begin{equation*}
-w_{i,i+\pmb{x}}=u_{i,i+\pmb{x}}=1+\tau^{3},\qquad
-w_{i,i+\pmb{y}}=u_{i,i+\pmb{y}}=(-)^{i_x}(1-\tau^{3}).
\end{equation*}
该拟设在平移 $\pmb{i}\to\pmb{i}+\pmb{x}$ 接规范变换 $G_x(\pmb{i})=(-)^{i_y}$ 下不变。它的 PSG 是式 (9.4.11) 中的 Z2B PSG。因此，$V<0$ 的基态是一个 Z2B 态。对 Z2B PSG，$G_xT_x$ 与 $G_yT_y$ 满足磁平移代数 $(G_xT_x)^{-1}(G_yT_y)^{-1}(G_xT_x)(G_yT_y)=-1$。费米子跳跃哈密顿量 $H_{\text{mean}}$ 描写费米子在磁场中的跳跃，每个方格有 $\pi$ 通量。不同的 PSG 告诉我们，$V<0$ 与 $V>0$ 的基态具有不同的量子序。

\section{正方格子上的 $Z_2$ 自旋液体与弦网凝聚}

\begin{keypoints}
\item 弦网凝聚导致演生的 $Z_2$ 规范理论和费米子。
\end{keypoints}

在这一节中，我们将从弦网凝聚的观点（Levin and Wen, 2003）出发，讨论精确可解自旋 $1/2$ 模型 (10.2.6)（Kitaev, 2003; Wen, 2003c）。该模型是最简单的模型之一，它演示了局域玻色模型中弦网凝聚与演生规范玻色子及费米子之间的联系。该模型也可以用从属玻色子方法求解，这使我们能够看到描写量子序的 PSG 如何与弦网凝聚相联系（见 10.4 节）。

\subsection{构造具有闭合弦网凝聚的哈密顿量}

\begin{keypoints}
\item 可以通过让某些弦状物体凝聚来构造具有演生规范场的局域玻色模型。
\end{keypoints}

让我们首先考虑正方格子上一个任意的自旋 $1/2$ 模型。我们想问的第一个问题是：什么样的自旋相互作用能够给出一个低能规范理论？如果我们相信规范理论与弦网理论之间的联系（Banks et al., 1977; Savit, 1980; Wen, 2003a），那么得到低能规范理论的一种途径，是设计一种自旋相互作用，使它允许大闭合弦的强烈涨落，却禁止其他类型的涨落（例如局域自旋翻转、开弦涨落等）。我们希望，大闭合弦的强烈涨落的出现将导致任意大小闭合弦的凝聚，继而给出一个低能规范理论。

\begin{figure}[H]
\centering
\includegraphics[width=0.72\textwidth]{fig_10.1.pdf}
\caption*{图 10.1\quad $H_J$ 的基态之上的一个开弦激发。}
\end{figure}

让我们从
\begin{equation*}
H_J=-J\sum_{\text{even}}\sigma_{\pmb{i}}^{x}-J\sum_{\text{odd}}\sigma_{\pmb{i}}^{y}
\end{equation*}
出发，其中 $\pmb{i}=(i_x,i_y)$ 标记格点，$\sigma^{x,y,z}$ 是泡利矩阵，而 $\sum_{\text{even}}$（或 $\sum_{\text{odd}}$）对偶格点求和，偶格点满足 $(-)^{\pmb{i}}\equiv(-1)^{i_x+i_y}=1$（或对奇格点求和，奇格点满足 $(-)^{\pmb{i}}=-1$）。$H_J$ 的基态，即 $|0\rangle$，具有这样的性质：偶格点上的自旋指向 $x$ 方向，奇格点上的自旋指向 $y$ 方向（见图 10.1）。这样一个态将被定义为无弦态（no-string state）。

为了产生一个弦激发，我们先画一条连接最近邻偶方格的弦（见图 10.1），然后把弦上的自旋翻转。这样一个弦态由下列弦产生算符（string creation operator，或简称弦算符）产生：
\begin{equation*}
W(C)=-\prod_{C}\sigma_{\pmb{i}}^{a_{\pmb{i}}}
\tag{10.3.1}
\end{equation*}
其中乘积 $\prod_{C}$ 遍及弦上的所有格点，且若 $\pmb{i}$ 为偶则 $a_{\pmb{i}}=y$，若 $\pmb{i}$ 为奇则 $a_{\pmb{i}}=x$。一个一般的弦态具有形式
\begin{equation*}
|C_1C_2...\rangle=W(C_1)W(C_2)...|0\rangle
\end{equation*}
其中 $C_1,C_2,\dots$ 是弦。这样的态将被称为弦网态（string-net state），因为弦可以相交和重叠。产生弦网的算符，即
\begin{equation*}
W(C_{\text{net}})=W(C_1)W(C_2)...
\end{equation*}
% （本块止于书页 452 末行显示公式；其后半句“将被称为弦网算符。”在书页 453 (p468) 顶部，由下一块收尾。）
弦网算符（string-net operator）。只要 $C_i$ 中至少有一个是开弦，态 $|C_1C_2\dots\rangle$ 就是开弦网态。相应的算符 $W(C_{\mathrm{net}})$ 将被称为开弦网算符。若所有 $C_i$ 都是闭合回路，则 $|C_1C_2\dots\rangle$ 就是闭合弦网态，而 $W(C_{\mathrm{net}})$ 是闭合弦网算符。

哈密顿量 $H_J$ 没有弦网凝聚，因为它的基态 $|0\rangle$ 不包含任何弦网，因而 $\langle0|W(C_{\mathrm{net}})|0\rangle=0$。要得到具有闭合弦网凝聚的哈密顿量，我们首先需要找到一个哈密顿量，其基态包含大量任意尺寸的闭合弦网，并且不包含任何开弦（或开弦网）。

让我们先写出一个使闭合弦和闭合弦网不耗费能量、而任何开弦都耗费大量能量的哈密顿量。这样一个哈密顿量可以取如下形式
\begin{equation*}
H_U=-U\sum_{\mathrm{even}}\hat F_{\pmb{i}}
\end{equation*}
其中
\begin{equation*}
\hat F_{\pmb{i}}=\sigma_{\pmb{i}}^{x}\sigma_{\pmb{i}+\pmb{x}}^{y}\sigma_{\pmb{i}+\pmb{x}+\pmb{y}}^{x}\sigma_{\pmb{i}+\pmb{y}}^{y}\tag{10.3.2}
\end{equation*}

我们发现，无弦态 $|0\rangle$ 是 $H_U$ 的基态之一（假定 $U>0$），能量为 $-UN_{\mathrm{site}}$。所有闭合弦态，例如 $W(C_{\mathrm{close}})|0\rangle$，也都是 $H_U$ 的基态，因为 $[H_U,W(C_{\mathrm{close}})]=0$。类似地，$[H_U,W(C_{\mathrm{net}})]=0$ 意味着任何闭合弦态也都是基态。

利用开弦算符 $W(C_{\mathrm{open}})$ 与 $\hat F_{\pmb{i}}$ 之间的对易关系
\begin{align*}
\hat F_{\pmb{i}}W(C_{\mathrm{open}})&=-W(C_{\mathrm{open}})\hat F_{\pmb{i}}, && \text{若 } C_{\mathrm{open}} \text{ 终止于方格 } \pmb{i},\\
\hat F_{\pmb{i}}W(C_{\mathrm{open}})&=+W(C_{\mathrm{open}})\hat F_{\pmb{i}}, && \text{其他情形},
\end{align*}
我们发现，开弦算符在其两个端点处翻转 $\hat F_{\pmb{i}}$ 的符号。开弦态 $W(C_{\mathrm{open}})|0\rangle$ 也是 $\hat F_{\pmb{i}}$ 的本征态，因而是 $H_U$ 的本征态，能量为 $-UN_{\mathrm{site}}+4U$。我们看到，开弦的每个端点耗费能量 $2U$。我们还注意到，闭合弦网的能量不依赖于网中弦的长度。因此，$H_U$ 中的弦没有张力。可以通过在哈密顿量中加入 $H_J$ 来引入弦张力，张力为每格点（或每段）$2J$。我们注意到，任何弦网态 $|C_1C_2\dots\rangle$ 都是 $H_U+H_J$ 的本征态。因此，由 $H_U+H_J$ 描述的弦网不涨落，从而不能凝聚。为使弦发生涨落，我们需要 $g$ 项
\begin{equation*}
H_g=-g\sum_{\pmb{p}}W(C_{\pmb{p}})=-g\sum_{\mathrm{odd}}\hat F_{\pmb{i}}
\end{equation*}
其中 $\pmb{p}$ 标记奇数方格，$C_{\pmb{p}}$ 是围绕方格 $\pmb{p}$ 的闭合弦。当 $H_g$ 作用在弦网态上时，它往弦网中添加一圈弦 $C_{\pmb{p}}$。因此 $g$ 项使弦发生涨落。我们的自旋 1/2 模型的总哈密顿量为
\begin{equation*}
H=H_U+H_J+H_g
\end{equation*}

当 $U\gg g\gg J$ 时，闭合弦几乎没有张力，$H_g$ 诱发的涨落不受张力的限制。因此，基态包含闭合弦的强烈涨落。换言之，基态被任意尺寸的闭合弦所充满。若 $U\gg g$，则闭合弦断开成开弦要耗费过多能量。因此，当 $U\gg g\gg J$ 时，我们自旋 1/2 模型的基态包含闭合弦凝聚（或者更确切地说，闭合弦网凝聚）。下面我们将研究这样一种弦网凝聚态的物理性质。

\subsection*{问题 10.3.1}

(a) 证明各 $\hat F_{\pmb{i}}$ 在 $x$ 与 $y$ 两个方向均为周期性边界条件的格子上满足
\begin{equation*}
\prod_{\pmb{i}}\hat F_{\pmb{i}}=1\tag{10.3.3}
\end{equation*}
式 (10.3.2) 中的负号使上述简单结果成为可能。

(b) 在两方向均为周期性边界条件的偶 $\times$ 偶格子上，证明各 $\hat F_{\pmb{i}}$ 满足更强的条件
\begin{equation*}
\prod_{\pmb{i}=\mathrm{even}}\hat F_{\pmb{i}}=1\tag{10.3.4}
\end{equation*}

\subsection{弦网凝聚与低能有效理论}

\begin{keypoints}\item 弦网凝聚给出规范激发。\end{keypoints}

让我们先讨论 $U\gg g>0$、$J=0$ 的弦网凝聚相。当 $J=0$ 时，模型 $H_U+H_g$ 成为上一节讨论过的模型 (10.2.6)，其中在偶数格点 $V_{\pmb{i}}=U$，在奇数格点 $V_{\pmb{i}}=g$。该模型是精确可解的，因为 $[\hat F_{\pmb{i}},\hat F_{\pmb{j}}]=0$（Kitaev, 2003）。$H_U+H_g$ 的本征态可以由 $\hat F_{\pmb{i}}$ 的共同本征态得到，即 $|\{F_{\pmb{i}}\}\rangle$，其中 $F_{\pmb{i}}$ 是 $\hat F_{\pmb{i}}$ 的本征值。由于 $\hat F_{\pmb{i}}^{2}=1$，$\hat F_{\pmb{i}}$ 的本征值就是 $F_{\pmb{i}}=\pm1$。$|\{F_{\pmb{i}}\}\rangle$ 的能量为 $E=-U\sum_{\pmb{i}=\mathrm{even}}F_{\pmb{i}}-g\sum_{\pmb{i}=\mathrm{odd}}F_{\pmb{i}}$。

注意到，对于尺寸为 $L_x\times L_y$ 的偶 $\times$ 偶周期格子上的有限系统，各 $\hat F_{\pmb{i}}$ 满足约束 $\prod_{\pmb{i}=\mathrm{even}}\hat F_{\pmb{i}}=1$ 与 $\prod_{\pmb{i}=\mathrm{odd}}\hat F_{\pmb{i}}=1$。因此，各 $F_{\pmb{i}}$ 并不独立，只能标记 $2^{L_xL_y}/4$ 种不同的组态。然而，我们自旋 1/2 模型的希尔伯特空间包含 $2^{L_xL_y}$ 个态。为了复现这 $2^{L_xL_y}$ 个态，$\hat F_{\pmb{i}}$ 的共同本征态必定是四重简并的。这与上一节得到的结果一致。

在 $U\gg g$ 的极限下，所有包含开弦的态都将具有量级为 $U$ 的能量。低能态只包含闭合弦网，并且满足
\begin{equation*}
F_{\pmb{i}}\big|_{\pmb{i}=\mathrm{even}}=1
\end{equation*}
不同的低能态由奇数方格上的 $F_{\pmb{i}}$ 标记，即 $F_{\pmb{i}}\big|_{\pmb{i}=\mathrm{odd}}=\pm1$。特别地，基态由下式给出
\begin{equation*}
F_{\pmb{i}}\big|_{\pmb{i}=\mathrm{odd}}=1.
\end{equation*}

这样的基态具有闭合弦凝聚。这是因为所有闭合弦算符 $W(C_{\mathrm{close}})$ 都与 $H_U+H_g$ 对易。于是，$H_U+H_g$ 的基态 $|\Psi_0\rangle$ 是 $W(C_{\mathrm{close}})$ 的本征态，并且满足
\begin{equation*}
\langle\Psi_0|W(C_{\mathrm{close}})|\Psi_0\rangle=1.
\end{equation*}

上式还意味着基态 $|\Psi_0\rangle$ 是各种闭合弦网组态的等权叠加。

基态之上的低能激发可以通过把某些奇数方格上的 $F_{\pmb{i}}$ 从 1 翻转到 $-1$ 得到。若把奇数方格上的 $F_{\pmb{i}}$ 视为 Z$_2$ 规范理论中的通量，我们发现模型的低能部分与 Z$_2$ 格点规范理论完全相同——至少对无限系统如此。这表明我们模型的低能有效理论是 Z$_2$ 格点规范理论。

然而，有人可能会反对这一结果，指出我们模型的低能部分也同样与在每个奇数方格上放一个自旋的伊辛模型相同，因此低能有效理论应当是伊辛模型。我们想指出，尽管对无限系统而言我们模型的低能部分与伊辛模型相同，但对有限系统而言，我们模型的低能部分不同于伊辛模型。例如，在具有周期性边界条件的有限偶 $\times$ 偶格子上，我们模型的基态有四重简并（Kitaev, 2003; Wen, 2003c）。伊辛模型没有这种简并，而 Z$_2$ 规范理论有这种简并。此外，我们的模型包含一种可以认定为 Z$_2$ 电荷（Z$_2$ charge）的激发。因此，我们模型的低能有效理论是 Z$_2$ 格点规范理论，而不是伊辛模型。奇数方格上 $F_{\pmb{i}}=-1$ 的激发可以视为 Z$_2$ 格点规范理论中的 Z$_2$ 涡旋激发。

为理解 Z$_2$ 电荷激发，我们注意到，在弦凝聚态中，弦不耗费能量，并且是不可观测的。把闭合弦算符作用到基态上，得到的仍是基态自身。由于弦不可观测，一段开弦的行为就像两个独立的粒子。开弦的每个端点对应偶数方格上的一个粒子。产生这样一个粒子所需的能量为 $U$ 的量级。现在，让我们考虑这样一个粒子围绕四个最近邻偶数方格的跳跃（见图 10.2）。每一步跳跃由一小段弦算符产生。我们看到，四个跳跃振幅的乘积由这四个偶数方格正中央那个奇数方格上 $\hat F_{\pmb{i}}$ 的本征值给出。这正是电荷与通量之间的关系。因此，若把奇数方格上的 $F_{\pmb{i}}$ 认定为 Z$_2$ 通量，则偶数方格上弦的端点就对应于 Z$_2$ 电荷。由于闭合弦凝聚，开弦的端点不被禁闭，彼此之间只有短程相互作用。因此，Z$_2$ 电荷的行为就像不挂着弦的点粒子。

\begin{figure}[H]
\centering
\includegraphics[width=0.72\textwidth]{fig_10.2.pdf}
\caption*{图 10.2\quad Z$_2$ 电荷围绕四个最近邻偶数方格的一次跳跃。}
\end{figure}

\subsection{三种类型的弦与演生费米子}

\begin{keypoints}\item 弦的端点是规范荷。某些弦的端点还可以是费米子。\end{keypoints}

下面我们将研究几种不同类型的弦。为避免混淆，我们把上面讨论的弦称为 T1 弦。T1 弦连接偶数方格（见图 10.3(a)）。

正如 Z$_2$ 电荷一样，一对 Z$_2$ 涡旋也由开弦算符产生。由于 Z$_2$ 涡旋对应于奇数方格上被翻转的 $\hat F_{\pmb{i}}$，产生 Z$_2$ 涡旋的开弦算符同样由式 (10.3.1) 给出，只是现在乘积遍及一条连接\emph{奇数}方格的弦。我们把这样的弦称为 T2 弦。

我们想指出，T2 弦的参考态（即无弦态）不同于 T1 弦的参考态。无 T2 弦态由 $|\tilde{0}\rangle$ 给出，其自旋在偶数格点指向 $y$ 方向、在奇数格点指向 $x$ 方向。由于 T1 弦与 T2 弦有不同的参考态，我们不能同时拥有 T1 弦和 T2 弦的稀薄气体。容易验证，T2 弦算符也与 $H_J+H_g$ 对易。因此，基态 $|\Psi_0\rangle$ 除了 T1 闭合弦的凝聚之外，还有 T2 闭合弦的凝聚。

\begin{figure}[H]
\centering
\includegraphics[width=0.72\textwidth]{fig_10.3.pdf}
\caption*{图 10.3\quad 三种类型的弦：(a) T1 弦；(b) T2 弦；(c) T3 弦。}
\end{figure}

Z$_2$ 涡旋的跳跃由一段短的 T2 开弦诱导。由于 T2 开弦算符彼此都对易，Z$_2$ 涡旋的行为像玻色子。类似地，Z$_2$ 电荷的行为也像玻色子。然而，T1 开弦算符与 T2 开弦算符不对易。结果是，T1 弦的端点与 T2 弦的端点具有非平凡的互统计。鉴于我们已经证明，让 Z$_2$ 电荷绕 Z$_2$ 涡旋运动一周会产生相位 $\pi$，Z$_2$ 电荷与 Z$_2$ 涡旋具有半子型（semionic）互统计。

T3 弦定义为 T1 弦与 T2 弦的束缚态（见图 10.3(c)）。T3 弦算符由一个 T1 弦算符与一个 T2 弦算符的乘积给出，具有如下形式
\begin{equation*}
W(C)=\prod_n\sigma_{\pmb{i}_n}^{l_n}
\end{equation*}
其中 $C$ 是一条连接相邻链（link）中点的弦，$\pmb{i}_n$ 是弦上的格点。这里，若弦在格点 $\pmb{i}_n$ 处不拐弯，则 $l_n=z$；若弦在格点 $\pmb{i}_n$ 处拐弯，则 $l_n=x$ 或 $y$；若拐弯构成右上角或左下角，则 $l_n=x$；若拐弯构成右下角或左上角，则 $l_n=y$（见图 10.3(c)）。基态也有 T3 闭合弦的凝聚。T3 弦的端点是由一条链两侧两个方格上的一个 Z$_2$ 电荷与一个 Z$_2$ 涡旋形成的束缚态（即链两侧的 $F_{\pmb{i}}=-1$）。这样的束缚态是费米子（见 7.1.2 节与图 7.8）。所以费米子生活在链上。有趣的是，我们模型中的弦网凝聚直接导致 Z$_2$ 规范结构和三种新类型的准粒子，即 Z$_2$ 电荷、Z$_2$ 涡旋和费米子。费米子作为开 T3 弦的端点，从我们这个纯玻色模型中演生出来！

由于 T1 弦的端点是 Z$_2$ 电荷，T1 弦可以视为 Z$_2$“电”通量的弦。类似地，T2 弦可以视为 Z$_2$“磁”通量的弦。

\section{用投影对称群对不同的弦网凝聚分类}

\subsection{四类弦网凝聚}

\begin{keypoints}\item 不同的弦网凝聚给出不同的量子有序相。\end{keypoints}

我们已经看到，当 $U>0$、$g>0$、$J=0$ 时，我们模型 $H_g+H_U+H_J$ 的基态由下式给出
\begin{equation*}
F_{\pmb{i}}\big|_{\pmb{i}=\mathrm{even}}=1,\qquad F_{\pmb{i}}\big|_{\pmb{i}=\mathrm{odd}}=1.
\end{equation*}
我们将把这样的相称为 Z$_2$ 相，以强调其低能 Z$_2$ 规范结构。在 Z$_2$ 相中，T1 弦算符 $W_1(C_1)$ 与 T2 弦算符 $W_2(C_2)$ 有如下期望值：
\begin{equation*}
\langle W_1(C_1)\rangle=1,\qquad \langle W_2(C_2)\rangle=1
\end{equation*}

当 $U>0$、$g<0$、$J=0$ 时，基态由下式给出
\begin{equation*}
F_{\pmb{i}}\big|_{\pmb{i}=\mathrm{even}}=1,\qquad F_{\pmb{i}}\big|_{\pmb{i}=\mathrm{odd}}=-1.
\end{equation*}
我们看到每个奇数方格都有 $\pi$ 通量。我们将把这样的相称为 Z$_2$ 通量相。T1 弦算符与 T2 弦算符有如下期望值：
\begin{equation*}
\langle W_1(C_1)\rangle=(-)^{N_{\mathrm{odd}}},\qquad \langle W_2(C_2)\rangle=1
\end{equation*}
其中 $N_{\mathrm{odd}}$ 是被 T1 弦 $C_1$ 包围的奇数方格的数目。

当 $U<0$、$g>0$、$J=0$ 时，基态由下式给出
\begin{equation*}
F_{\pmb{i}}\big|_{\pmb{i}=\mathrm{even}}=-1,\qquad F_{\pmb{i}}\big|_{\pmb{i}=\mathrm{odd}}=1.
\end{equation*}
每个偶数方格上有一个 Z$_2$ 电荷。我们将把这样的相称为 Z$_2$ 电荷相。T1 弦算符与 T2 弦算符有如下期望值：
\begin{equation*}
\langle W_1(C_1)\rangle=1,\qquad \langle W_2(C_2)\rangle=(-)^{N_{\mathrm{even}}}
\end{equation*}
其中 $N_{\mathrm{even}}$ 是被 T2 弦 $C_2$ 包围的偶数方格的数目。注意，Z$_2$ 通量相与 Z$_2$ 电荷相仅相差一个格子平移，本质上是同一个相。

当 $U<0$、$g<0$、$J=0$ 时，基态变为
\begin{equation*}
F_{\pmb{i}}\big|_{\pmb{i}=\mathrm{even}}=-1,\qquad F_{\pmb{i}}\big|_{\pmb{i}=\mathrm{odd}}=-1.
\end{equation*}
每个偶数方格上有一个 Z$_2$ 电荷，且每个奇数方格有 $\pi$ 通量。我们将把这样的相称为 Z$_2$ 通量--电荷相。T1 弦算符与 T2 弦算符有如下期望值：
\begin{equation*}
\langle W_1(C_1)\rangle=(-)^{N_{\mathrm{odd}}},\qquad \langle W_2(C_2)\rangle=(-)^{N_{\mathrm{even}}}
\end{equation*}

当 $U=g=0$、$J\neq0$ 时，模型同样精确可解。基态是一个简单的自旋极化态（见图 10.1）。

从 $H_U+H_g+H_J$ 模型的这些精确可解极限出发，我们提出一幅相图，草图示于图 10.4。该相图包含四个不同的弦网凝聚相和一个无弦凝聚的相。所有的相都有相同的对称性，仅凭它们不同的量子序相互区分。

\subsection{投影对称群与凝聚弦的端点}

\begin{keypoints}\item PSG 所描述的投影对称性，正是凝聚弦端点有效理论的对称性。\end{keypoints}

从各不相同的 $\langle W_1(C_1)\rangle$ 与 $\langle W_2(C_2)\rangle$ 我们看到，前四个相具有不同的弦网凝聚。然而，它们都有相同的对称性。这就引出一个问题：在没有对称性破缺的情况下，我们怎么知道上述四个相的确是不同的相？我们怎么知道不可能在不经历相变的情况下把一个弦网凝聚态变成另一个弦网凝聚态？

\begin{figure}[H]
\centering
\includegraphics[width=0.72\textwidth]{fig_10.4.pdf}
\caption*{图 10.4\quad $H=H_U+H_g+H_J$ 模型的建议相图。这里假定 $J$ 为正。四个弦网凝聚相由一对 PSG（$PSG_{\mathrm{charge}}$，$PSG_{\mathrm{vortex}}$）表征。SP 标记自旋极化相。}
\end{figure}

下面我们将证明，不同的弦网凝聚可以用不同的 PSG 来描述（正如不同的对称性破缺序可以用基态的不同对称群来描述）。在第 9 章中，不同的量子序是通过它们不同的 PSG 引入的（Wen, 2002b,c）。那里证明了 PSG 是量子相的普适性质，只有相变才能改变 PSG。因此，不同的弦网凝聚态由不同的 PSG 表征，这一事实表明这些不同的弦网凝聚态属于不同的量子相。

弦网凝聚与 PSG 之间的联系也使我们能够把弦网凝聚与第 9 章引入的量子序联系起来。事实上，第 9 章讨论的投影构造可以视为构造弦网凝聚态的一种方法。

为看清 PSG 与弦网凝聚之间的联系，我们注意到，当闭合弦网凝聚时，开弦的端点表现得像独立的粒子。让我们考虑描述 T1 弦两个端点的双粒子态 $|\pmb{p}_1\pmb{p}_2\rangle$。注意，T1 弦的端点（因而 Z$_2$ 电荷）只生活在偶数方格上。因此，$\pmb{p}_1$ 与 $\pmb{p}_2$ 只标记偶数方格。对于我们的模型 $H_U+H_g$，$|\pmb{p}_1\pmb{p}_2\rangle$ 是能量本征态，而且 Z$_2$ 电荷不跳跃。这里我们想往哈密顿量中加入一项
\begin{equation*}
H_t=t\sum_{\pmb{i}}(\sigma_{\pmb{i}}^{x}+\sigma_{\pmb{i}}^{y})+t'\sum_{\pmb{i}}\sigma_{\pmb{i}}^{z}
\end{equation*}
$t$ 项 $t\sum_{\pmb{i}}(\sigma_{\pmb{i}}^{x}+\sigma_{\pmb{i}}^{y})$ 使 Z$_2$ 电荷在偶数方格之间跳跃，跳跃幅为 $t$ 的量级。两个 Z$_2$ 电荷的动力学由双粒子希尔伯特空间中的如下有效哈密顿量描述：
\begin{equation*}
H=H(\pmb{p}_1)+H(\pmb{p}_2)
\end{equation*}
其中 $H(\pmb{p}_1)$ 描述第一个粒子 $\pmb{p}_1$ 的跳跃，$H(\pmb{p}_2)$ 描述第二个粒子 $\pmb{p}_2$ 的跳跃。现在我们可以定义弦网凝聚态的 PSG 了：PSG 就是跳跃哈密顿量 $H(\pmb{p})$ 的对称群。

由于底层模型 $H_U+H_g+H_t$ 具有平移对称性，我们可能天真地预期 Z$_2$ 电荷的跳跃哈密顿量 $H(\pmb{p})$ 在 $\pmb{x}+\pmb{y}$ 与 $\pmb{x}-\pmb{y}$ 方向也具有平移对称性：
\begin{align*}
H(\pmb{p})&=T_{xy}^{\dagger}H(\pmb{p})T_{xy}, & T_{xy}|\pmb{p}\rangle&=|\pmb{p}+\pmb{x}+\pmb{y}\rangle\\
H(\pmb{p})&=T_{x\bar{y}}^{\dagger}H(\pmb{p})T_{x\bar{y}}, & T_{x\bar{y}}|\pmb{p}\rangle&=|\pmb{p}+\pmb{x}-\pmb{y}\rangle\tag{10.4.1}
\end{align*}
如果上式成立，则意味着 PSG 与平移对称群相同。

结果表明，式 (10.4.1) 过强。即使 $H(\pmb{p})$ 不满足式 (10.4.1)，底层自旋模型仍然可以具有平移对称性。然而，$H(\pmb{p})$ 可能的对称群（即各个 PSG）受到底层自旋模型平移对称性的很强约束。下面我们将解释，为什么 PSG 可以不同于物理自旋模型的对称群，以及为与底层自旋模型的平移对称性相一致，PSG 必须满足什么条件。

我们注意到，弦总有两个端点。因此，一个物理态总是拥有偶数个 Z$_2$ 电荷。平移作用在双粒子态上为
\begin{align*}
T_{xy}^{(2)}|\pmb{p}_1,\pmb{p}_2\rangle&=|\pmb{p}_1+\pmb{x}+\pmb{y},\ \pmb{p}_2+\pmb{x}+\pmb{y}\rangle\\
T_{x\bar{y}}^{(2)}|\pmb{p}_1,\pmb{p}_2\rangle&=|\pmb{p}_1+\pmb{x}-\pmb{y},\ \pmb{p}_2+\pmb{x}-\pmb{y}\rangle
\end{align*}
这里 $T_{xy}^{(2)}$ 与 $T_{x\bar{y}}^{(2)}$ 满足平移代数
\begin{equation*}
T_{xy}^{(2)}T_{x\bar{y}}^{(2)}=T_{x\bar{y}}^{(2)}T_{xy}^{(2)}\tag{10.4.2}
\end{equation*}

我们注意到，$T_{xy}^{(2)}$ 与 $T_{x\bar{y}}^{(2)}$ 是单粒子态上平移算符的直积。因此，在某种意义上，单粒子平移是双粒子平移的平方根。双粒子平移代数对单粒子平移代数施加了很强的约束。

单粒子平移最一般的形式由 $T_{xy}G_{xy}$ 与 $T_{x\bar{y}}G_{x\bar{y}}$ 给出，其中算符 $T_{xy,x\bar{y}}$ 与 $G_{xy,x\bar{y}}$ 的作用定义为
\begin{align*}
T_{xy}|\pmb{p}\rangle&=|\pmb{p}+\pmb{x}+\pmb{y}\rangle, & T_{x\bar{y}}|\pmb{p}\rangle&=|\pmb{p}+\pmb{x}-\pmb{y}\rangle,\\
G_{xy}|\pmb{p}\rangle&=\mathrm{e}^{\phi_{xy}(\pmb{p})}|\pmb{p}\rangle, & G_{x\bar{y}}|\pmb{p}\rangle&=\mathrm{e}^{\phi_{x\bar{y}}(\pmb{p})}|\pmb{p}\rangle.
\end{align*}
为了使直积 $T_{xy}^{(2)}=T_{xy}G_{xy}\otimes T_{xy}G_{xy}$ 与 $T_{x\bar{y}}^{(2)}=T_{x\bar{y}}G_{x\bar{y}}\otimes T_{x\bar{y}}G_{x\bar{y}}$ 复现平移代数 (10.4.2)，我们只要求 $T_{xy}G_{xy}$ 与 $T_{x\bar{y}}G_{x\bar{y}}$ 满足
\begin{equation*}
T_{xy}G_{xy}T_{x\bar{y}}G_{x\bar{y}}=T_{x\bar{y}}G_{x\bar{y}}T_{xy}G_{xy}\tag{10.4.3}
\end{equation*}
或
\begin{equation*}
T_{xy}G_{xy}T_{x\bar{y}}G_{x\bar{y}}=-T_{x\bar{y}}G_{x\bar{y}}T_{xy}G_{xy}\tag{10.4.4}
\end{equation*}

算符 $T_{xy}G_{xy}$ 与 $T_{x\bar{y}}G_{x\bar{y}}$ 生成一个群，这个群正是 PSG。两个不同的代数 (10.4.3) 与 (10.4.4) 生成两个不同的 PSG，二者都与作用在双粒子态上的平移群相一致。我们把由式 (10.4.3) 生成的 PSG 称为 Z2a PSG，把由式 (10.4.4) 生成的 PSG 称为 Z2b PSG。

让我们回顾一下 PSG 的一般定义。PSG 是一个群，它是对称群（SG）的扩张，即 PSG 包含一个正规子群（称为不变规范群或 IGG），使得
\begin{equation*}
PSG/IGG=SG
\end{equation*}

在我们的情形中，SG 是平移群 $SG=\{1,T_{xy}^{(2)},T_{x\bar{y}}^{(2)},\dots\}$。IGG 由单粒子态上满足 $G_0\otimes G_0=1$ 的变换 $G_0$ 构成。我们发现 IGG 由下式生成：
\begin{equation*}
G_0|\pmb{p}\rangle=-|\pmb{p}\rangle
\end{equation*}

单粒子态上的这三个变换 $G_0$、$T_{xy}G_{xy}$ 与 $T_{x\bar{y}}G_{x\bar{y}}$ 生成 Z2a 或 Z2b PSG。

我们现在看到，底层平移对称性并不要求单粒子跳跃哈密顿量 $H(\pmb{p})$ 具有平移对称性，只要求 $H(\pmb{p})$ 在 Z2a PSG 或 Z2b PSG 下不变。当 $H(\pmb{p})$ 在 Z2a PSG 下不变时，跳跃哈密顿量具有通常的平移对称性。当 $H(\pmb{p})$ 在 Z2b PSG 下不变时，跳跃哈密顿量具有磁平移对称性，描述磁场中每个奇数方格带 $\pi$ 通量的跳跃。

\subsection{投影对称群对不同的弦网凝聚的分类}

\begin{keypoints}\item 不同的弦网凝聚有不同的 PSG。由于 PSG 是量子相的普适性质，我们发现不同的弦网凝聚对应于不同的相。\end{keypoints}

在了解了弦端点的跳跃哈密顿量可能具有的 PSG 之后，我们现在可以计算实际的 PSG 了。让我们考虑模型 $H_U+H_g+H_t$ 的两个基态：一个具有 $F_{\pmb{i}}\big|_{\pmb{i}=\mathrm{odd}}=1$（对应 $g>0$），另一个具有 $F_{\pmb{i}}\big|_{\pmb{i}=\mathrm{odd}}=-1$（对应 $g<0$）。两个基态在 $\pmb{x}+\pmb{y}$ 与 $\pmb{x}-\pmb{y}$ 方向都有相同的平移对称性。然而，相应的单粒子跳跃哈密顿量 $H(\pmb{p})$ 却有不同的对称性。对于 $F_{\pmb{i}}\big|_{\pmb{i}=\mathrm{odd}}=1$ 的态，奇数方格没有通量，$H(\pmb{p})$ 具有通常的平移对称性，在 Z2a PSG 下不变。对于 $F_{\pmb{i}}\big|_{\pmb{i}=\mathrm{odd}}=-1$ 的态（译注：原文此处印作“$F_{\pmb{i}}|_{\pmb{i}=\mathrm{odd}}=1$”，按上下文应为 $-1$），奇数方格有 $\pi$ 通量，$H(\pmb{p})$ 具有磁平移对称性，其 PSG 为 Z2b PSG。因此，$F_{\pmb{i}}\big|_{\pmb{i}=\mathrm{odd}}=1$ 的态与 $F_{\pmb{i}}\big|_{\pmb{i}=\mathrm{odd}}=-1$ 的态尽管对称性相同，却有不同的序。两个态中不同的量子序可以用它们不同的 PSG 来表征。

如 10.4.1 小节所述，上述两个态有不同的弦网凝聚。对 $\hat F_{\pmb{i}}\big|_{\pmb{i}=\mathrm{odd}}=1$ 的态，T1 闭合弦算符的平均值为 $\langle W(C)\rangle=1$；而对 $\hat F_{\pmb{i}}\big|_{\pmb{i}=\mathrm{odd}}=-1$ 的态，T1 闭合弦算符的平均值为 $\langle W(C)\rangle=(-)^{N_{\mathrm{odd}}(C)}$，其中 $N_{\mathrm{odd}}(C)$ 是被 T1 闭合弦 $C$ 包围的奇数方格的数目。因此，我们也可以说，不同的弦网凝聚可以用不同的 PSG 来表征。

在上面的讨论中，我们只在一些特定的精确可解模型中表明了 PSG 可以表征不同的弦网凝聚态。事实上，在一般模型中，不同的 PSG 确实表征不同的量子相。这是因为，正如 9.9.1 小节所证明的，PSG 是相的普适性质（Wen, 2002c）。弦网凝聚态的 PSG 对物理模型的任何不改变模型对称性的局域微扰都是稳健的。

上述讨论同样适用于 Z$_2$ 涡旋与 T2 弦。因此，我们模型中的量子序由一对 PSG（$PSG_{\mathrm{charge}}$，$PSG_{\mathrm{vortex}}$）描述，一个对应 Z$_2$ 电荷，一个对应 Z$_2$ 涡旋。PSG 对（$PSG_{\mathrm{charge}}$，$PSG_{\mathrm{vortex}}$）使我们能够区分模型 $H=H_U+H_g+H_t$ 的四个不同的弦网凝聚态（见图 10.5）。

\subsection{T3 弦端点的投影对称群}

\begin{keypoints}\item 对于具有几种不同类型弦凝聚的态，不同凝聚弦的端点可以有不同的投影对称性和不同的 PSG。\end{keypoints}

\begin{figure}[H]
\centering
\includegraphics[width=0.72\textwidth]{fig_10.5.pdf}
\caption*{图 10.5\quad $H=H_U+H_g+H_t$ 模型的建议相图。这里假定 $t=t'$ 为正。四个弦网凝聚相由 PSG 对（$PSG_{\mathrm{charge}}$，$PSG_{\mathrm{vortex}}$）表征。SP 标记均匀自旋极化相。所有的相都有相同的对称性，并且没有任何相变改变对称性。}
\end{figure}

\begin{keypoints}\item 9.4.2 小节引入的 PSG 对应于 T3 弦的 PSG。\end{keypoints}

在本小节中，我们将假设模型中 $U=g=V$：
\begin{equation*}
H_U+H_g+H_t=H_t-V\sum_{\pmb{i}}\hat F_{\pmb{i}}\tag{10.4.5}
\end{equation*}

新模型具有由 $\pmb{i}\to\pmb{i}+\pmb{x}$ 与 $\pmb{i}\to\pmb{i}+\pmb{y}$ 生成的更大的平移对称性（见图 10.5）。注意到，当 $H_t=0$ 时（译注：原文此处印作“$H_l=0$”，应为 $H_t$），新模型与 10.2.1 小节研究过的自旋哈密顿量 (10.2.6) 相同。凝聚的 T1 弦与 T2 弦的 PSG 已在上一小节研究过，这里我们想讨论 T3 弦的 PSG。由于 T3 弦的端点生活在链上，相应的单粒子跳跃哈密顿量 $H_f(\pmb{l})$ 描述费米子在链之间的跳跃。显然，$H_f(\pmb{l})$ 的对称群（即 PSG）可以不同于 $H(\pmb{p})$ 的对称群。

让我们把费米子跳跃 $\pmb{l}\to\pmb{l}+\pmb{x}$ 定义为两步跳跃 $\pmb{l}\to\pmb{l}+\frac{\pmb{x}}{2}-\frac{\pmb{y}}{2}\to\pmb{l}+\pmb{x}$ 的组合（见图 10.6）。跳跃 $\pmb{l}\to\pmb{l}+\frac{\pmb{x}}{2}-\frac{\pmb{y}}{2}$ 由 $\sigma_{\pmb{l}+\frac{\pmb{x}}{2}}^{x}$ 产生，跳跃 $\pmb{l}+\frac{\pmb{x}}{2}-\frac{\pmb{y}}{2}\to\pmb{l}+\pmb{x}$ 由 $\sigma_{\pmb{l}+\frac{\pmb{x}}{2}}^{y}$ 产生。因此，跳跃 $\pmb{l}\to\pmb{l}+\pmb{x}$ 由 $\sigma_{\pmb{l}+\frac{\pmb{x}}{2}}^{y}\sigma_{\pmb{l}+\frac{\pmb{x}}{2}}^{x}$ 产生。类似地，我们把费米子跳跃 $\pmb{l}\to\pmb{l}+\pmb{y}$ 定义为两步跳跃 $\pmb{l}\to\pmb{l}+\frac{\pmb{x}}{2}+\frac{\pmb{y}}{2}\to\pmb{l}+\pmb{y}$，它由 $\sigma_{\pmb{l}+\frac{\pmb{x}}{2}+\pmb{y}}^{x}\sigma_{\pmb{l}+\frac{\pmb{x}}{2}}^{y}$ 产生。
\begin{figure}[H]
\centering
\includegraphics[width=0.72\textwidth]{fig_10.6.pdf}
\caption*{图 10.6\quad 费米子绕一个方格和绕一个格点的跳跃。}
\end{figure}
在这样的定义下，费米子绕方格 $\pmb{l}\to\pmb{l}+\pmb{x}\to\pmb{l}+\pmb{x}+\pmb{y}\to\pmb{l}+\pmb{y}\to\pmb{l}$ 跳跃一周由下式产生（见图 10.6）：$(\sigma_1^{y}\sigma_4^{x})(\sigma_4^{x}\sigma_4^{y})(\sigma_3^{x}\sigma_2^{y})(\sigma_1^{y}\sigma_1^{x})=\sigma_4^{y}\sigma_3^{x}\sigma_2^{y}\sigma_1^{x}=\hat F_1$，或者 $(\sigma_5^{z})(\sigma_5^{y}\sigma_8^{x})(\sigma^{z})(\sigma_7^{x}\sigma_6^{y})=-\sigma_5^{x}\sigma_8^{y}\sigma_7^{x}\sigma_6^{y}=-\hat F_5$。基态具有 $F_{\pmb{i}}=\mathrm{sgn}(V)$。然而，由于 T3 弦终止在链 $5\to6$ 上，我们有 $-\hat F_5=\hat F_{\pmb{i}}=\mathrm{sgn}(V)$；另一方面，我们有 $\hat F_1=\hat F_{\pmb{i}}=\mathrm{sgn}(V)$。因此，费米子绕 1234 与 5678 两个方格跳跃的总振幅由式 (10.4.5) 中 $V$ 的符号给出。当 $\mathrm{sgn}(V)=1$ 时，费米子看不到通量；而当 $\mathrm{sgn}(V)=-1$ 时，费米子看到每个方格有 $\pi$ 通量。结果是，$H_f(\pmb{l})$ 平移对称性的生成元 $(T_xG_x,\,T_yG_y)$ 满足磁平移代数（见式 (7.2.3)）
\begin{equation*}
(T_yG_y)^{-1}(T_xG_x)^{-1}T_yG_yT_xG_x=\mathrm{sgn}(V).\tag{10.4.6}
\end{equation*}

除了在 $(T_xG_x,\,T_yG_y)$ 下不变之外，$H_f(\pmb{l})$ 在 $G_0$ 下也不变：
\begin{equation*}
G_0|\pmb{l}\rangle=-|\pmb{l}\rangle
\end{equation*}
于是，$(G_0,\,T_xG_x,\,T_yG_y)$ 生成 $H_f(\pmb{l})$ 的对称群——费米子 PSG。

当 $\mathrm{sgn}(V)=1$ 时，$T_xG_x$ 与 $T_yG_y$ 满足平移代数，相应的费米子 PSG 是 Z2A PSG。当 $\mathrm{sgn}(V)=-1$ 时，$T_xG_x$ 与 $T_yG_y$ 满足每个方格带 $\pi$ 通量的磁平移代数，相应的费米子 PSG 是 Z2B PSG。我们看到，基态的量子序也可以用费米子 PSG 来表征。$V<0$ 基态的量子序由 Z2A PSG 表征，而 $V>0$ 基态的量子序由 Z2B PSG 表征。

在 10.2 节中，自旋 1/2 模型 (10.4.5)（取 $t=t'=0$）曾通过把自旋拆成两个费米子、用从属玻色子方法求解。那里证明了，$V>0$ 与 $V<0$ 两种态的费米子跳跃哈密顿量具有不同的对称性，或者说在不同的 PSG 下不变。不同的 PSG 意味着 $V>0$ 与 $V<0$ 基态中不同的量子序。10.2 节中针对 $V>0$ 与 $V<0$ 相得到的 PSG 与
这里我们想指出，9.4 节中引入的 PSG 都是费米子 PSG。它们只是可用于刻画量子序的众多不同种类的 PSG 中的一种。一般地说，一个量子有序态可以包含几种不同类型弦的凝聚，每一类凝聚弦的端点都将拥有各自的 PSG。找出所有类型的凝聚弦及其 PSG，将使我们能够对量子序作出更完备的刻画。

\section{立方格子上的演生费米子与弦网凝聚}

\begin{keypoints}
\item 三维模型中的演生费米子，无法用把通量附着到玻色子上的方式来理解。要理解三维中的演生费米子，确实需要弦网理论。
\end{keypoints}

在 $2+1$ 维中，把 $\pi$ 通量束缚到一个单位电荷上（或者把一个 $Z_2$ 涡旋束缚到一个 $Z_2$ 电荷上），便可以得到一个费米子。在我们的自旋 $1/2$ 模型 $H_U+H_g$ 中，我们正是这样得到费米子的。由于 $Z_2$ 涡旋与 $Z_2$ 电荷都以开弦端点的形式出现，费米子也以弦端点的形式出现。我们看到，对于 $(2+1)$ 维玻色模型中的演生费米子，我们有两套理论：演生费米子既可以视为电荷与通量的束缚态，也可以视为弦的端点。然而，在 $3+1$ 维中，我们不能通过附着 $\pi$ 通量把玻色子变成费米子。于是人们或许会问：费米子能否从 $(3+1)$ 维的局域玻色模型中演生出来？本节我们将研究立方格子上一个精确可解的自旋 $3/2$ 模型。我们将研究这样一个模型中的弦网凝聚态，并证明费米子仍然可以作为弦的端点出现。事实上，关于演生费米子的弦网理论在任意维数下都成立。

\subsection{立方格子上的精确可解自旋 $3/2$ 模型}

为了在三维立方格子上构造一个精确可解模型，我们先引入算符 $\gamma_{\pmb{i}}^{ab}$，$a,b=x,\bar{x},y,\bar{y},z,\bar{z}$，它们作用在每个格点的局域希尔伯特空间之内。在一个格点上，$\gamma^{ab}$ 满足以下代数（格点指标 $\pmb{i}$ 已略去）：
\begin{align*}
\gamma^{ab}&=-\gamma^{ba}=(\gamma^{ab})^{\dagger}, & (\gamma^{ab})^{2}&=1, \tag{10.5.1}\\
[\gamma^{ab},\gamma^{cd}]&=0, & \gamma^{ab}\gamma^{bc}&=\mathrm{i}\gamma^{ac}, \qquad \text{若 } a,b,c,d \text{ 互不相同}
\end{align*}

为了证明这个代数是自洽的，令 $\lambda^{a}$ 为一个由 $a=x,\bar{x},y,\bar{y},z,\bar{z}$ 标记的马约拉纳费米子算符。从马约拉纳费米子的代数 $\{\lambda^{a},\lambda^{b}\}=2\delta_{ab}$ 出发可以证明，
\begin{equation*}
\gamma^{ab}=\frac{\mathrm{i}}{2}\left(\lambda_{\pmb{i}}^{a}\lambda_{\pmb{i}}^{b}-\lambda_{\pmb{i}}^{b}\lambda_{\pmb{i}}^{a}\right)
\end{equation*}
满足上述代数。

上述代数有一个四维表示。我们首先注意到，$\gamma^{az}\equiv\gamma^{a}$（$a=x,\bar{x},y,\bar{y}$）满足狄拉克代数
\begin{equation*}
\{\gamma^{a},\gamma^{b}\}=2\delta_{ab}, \qquad a,b=x,\bar{x},y,\bar{y}. \tag{10.5.2}
\end{equation*}
因此，我们可以用四个 $4\times4$ 狄拉克矩阵把它们表示为
\begin{align*}
\gamma^{xz}=\gamma^{x}=\sigma^{x}\otimes\sigma^{x}, & & \gamma^{\bar{x}z}=\gamma^{\bar{x}}=\sigma^{y}\otimes\sigma^{x}, &\\
\gamma^{yz}=\gamma^{y}=\sigma^{z}\otimes\sigma^{x}, & & \gamma^{\bar{y}z}=\gamma^{\bar{y}}=\sigma^{0}\otimes\sigma^{y}. & \tag{10.5.3}
\end{align*}
我们还可以定义
\begin{equation*}
\gamma^{5}=\gamma^{x}\gamma^{\bar{x}}\gamma^{y}\gamma^{\bar{y}}=-\sigma^{0}\otimes\sigma^{z}
\end{equation*}
它与 $\gamma^{a}$ 反对易，即 $\{\gamma^{a},\gamma^{5}\}=0$。由于 $\gamma^{z\bar{z}}$ 与 $\gamma^{az}\equiv\gamma^{a}$（$a=x,\bar{x},y,\bar{y}$）反对易，我们可以把 $\gamma^{5}$ 认同为 $\gamma^{z\bar{z}}$。利用代数 (10.5.1)，可以把其余的 $\gamma^{ab}$ 表示为
\begin{equation*}
\gamma^{a\bar{z}}=\mathrm{i}\gamma^{a}\gamma^{5}, \qquad a=x,\bar{x},y,\bar{y}
\end{equation*}
这样，我们就把全部 $\gamma^{ab}$（$a,b=x,\bar{x},y,\bar{y},z,\bar{z}$）用狄拉克矩阵 $\gamma^{x,\bar{x},y,\bar{y}}$ 表示了出来。可以证明，这样定义的 $\gamma^{ab}$ 满足代数 (10.5.1)。

由于 $\gamma^{ab}$ 是四维的，我们可以说它们作用在自旋 $3/2$ 态上。借助 $\gamma_{\pmb{i}}^{ab}$，我们可以在立方格子上写出一个精确可解的自旋 $3/2$ 模型如下：
\begin{align*}
H_{3D}&\\
&=g\sum_{\pmb{i}}\Big(\gamma_{\pmb{i}}^{yx}\gamma_{\pmb{i}+\pmb{x}}^{\bar{x}y}\gamma_{\pmb{i}+\pmb{x}+\pmb{y}}^{\bar{y}\bar{x}}\gamma_{\pmb{i}+\pmb{y}}^{x\bar{y}}+\gamma_{\pmb{i}}^{zy}\gamma_{\pmb{i}+\pmb{y}}^{\bar{y}z}\gamma_{\pmb{i}+\pmb{y}+\pmb{z}}^{\bar{z}\bar{y}}\gamma_{\pmb{i}+\pmb{z}}^{y\bar{z}}+\gamma_{\pmb{i}}^{xz}\gamma_{\pmb{i}+\pmb{z}}^{\bar{z}x}\gamma_{\pmb{i}+\pmb{z}+\pmb{x}}^{\bar{x}\bar{z}}\gamma_{\pmb{i}+\pmb{x}}^{z\bar{y}}\Big)\\
&=-g\sum_{\pmb{p}}\hat F_{\pmb{p}} \tag{10.5.4}
\end{align*}
其中 $\pmb{p}$ 标记立方格子中的方格，而视方格 $\pmb{p}$ 的取向，$\hat F_{\pmb{p}}$ 等于 $-\gamma_{\pmb{i}}^{yx}\allowbreak\gamma_{\pmb{i}+\pmb{x}}^{\bar{x}y}\allowbreak\gamma_{\pmb{i}+\pmb{x}+\pmb{y}}^{\bar{y}\bar{x}}\allowbreak\gamma_{\pmb{i}+\pmb{y}}^{x\bar{y}}$、$-\gamma_{\pmb{i}}^{zy}\allowbreak\gamma_{\pmb{i}+\pmb{y}}^{\bar{y}z}\allowbreak\gamma_{\pmb{i}+\pmb{y}+\pmb{z}}^{\bar{z}\bar{y}}\allowbreak\gamma_{\pmb{i}+\pmb{z}}^{y\bar{z}}$ 或 $-\gamma_{\pmb{i}}^{xz}\allowbreak\gamma_{\pmb{i}+\pmb{z}}^{\bar{z}x}\allowbreak\gamma_{\pmb{i}+\pmb{z}+\pmb{x}}^{\bar{x}\bar{z}}\allowbreak\gamma_{\pmb{i}+\pmb{x}}^{z\bar{y}}$。

注意到
\begin{equation*}
[\hat F_{\pmb{p}},\hat F_{\pmb{p}}']=0, \qquad (\hat F_{\pmb{p}})^{2}=1,
\end{equation*}
就容易求出 $H_{3D}$ 的本征值。设 $|\{F_{\pmb{p}}\},\alpha\rangle$ 是 $\hat F_{\pmb{p}}$ 的共同本征态，即 $\hat F_{\pmb{p}}|\{F_{\pmb{p}}\},\alpha\rangle=F_{\pmb{p}}|\{F_{\pmb{p}}\},\alpha\rangle$，其中 $\alpha$ 标记可能的简并。那么 $|\{F_{\pmb{p}}\},\alpha\rangle$ 也是能量本征态，能量为 $-g\sum_{\pmb{p}}F_{\pmb{p}}$。由于 $F_{\pmb{p}}=\pm1$，在周期性格子上基态能量为 $E_0=-3|g|N_{\text{site}}$，其中 $N_{\text{site}}$ 是格点数目。激发态可以通过翻转某些 $F_{\pmb{p}}$ 的符号得到，其能量为 $E_0$ 加上 $2|g|$ 的整数倍。

有一个微妙之处：诸 $F_{\pmb{p}}$ 并不全部独立。若 $S$ 是一个单胞立方体的表面，则有算符恒等式
\begin{equation*}
\prod_{\pmb{p}\in S}\hat F_{\pmb{p}}=1
\end{equation*}
由此可知，对所有立方体都有 $\prod_{\pmb{p}\in S}F_{\pmb{p}}=1$。如果我们把 $F_{\pmb{p}}$ 看作穿过方格 $\pmb{p}$ 的 $Z_2$ 通量，那么上述约束就意味着通量守恒。这意味着我们模型的能谱与立方格子上一个 $Z_2$ 规范理论的能谱相同。基态可以看作无通量的态，即对所有 $\pmb{p}$ 都有 $F_{\pmb{p}}=1$。基态之上的激发由通量圈描述。基本激发对应于最小的通量圈，即在与某条链 $\langle\pmb{i}\pmb{j}\rangle$ 相邻的四个方格 $\pmb{p}$ 上 $F_{\pmb{p}}=-1$。我们可以把这些激发看作生活在立方格子\emph{链}上的准粒子（见图 10.7）。正如我们后面将看到的，这种通量圈激发具有费米统计。

\subsection*{问题 10.5.1}

\begin{enumerate}
\item 证明式 (10.5.3) 中的 $4\times4$ 矩阵 $\gamma^{a}$ 满足狄拉克代数。
\item 证明式 (10.5.4) 求和中的各项彼此对易。（提示：可以利用 $\gamma^{ab}$ 的马约拉纳费米子表示。）这使得我们能够精确求解 $H_{3D}$。
\end{enumerate}

\subsection*{问题 10.5.2}

\textbf{用马约拉纳费米子和投影构造求解式 (10.5.4)}

\begin{enumerate}
\item 设 $\lambda_{\pmb{i}}^{a}$ 是马约拉纳费米子算符，其中 $a=x,\bar{x},y,\bar{y},z,\bar{z}$。对每条最近邻链，定义
\begin{equation*}
\hat U_{\pmb{i}\pmb{j}}=\lambda_{\pmb{i}}^{a}\lambda_{\pmb{j}}^{b} \tag{10.5.5}
\end{equation*}
其中 $ab$ 是与链 $\pmb{i}\to\pmb{j}$ 相联系的指标对（见图 10.7）。证明诸 $\hat U_{\pmb{i}\pmb{j}}$ 构成一组相互对易的算符。
\item 考虑费米子哈密顿量
\begin{equation*}
H=-\sum_{\pmb{p}}g\hat F_{\pmb{p}}, \qquad \hat F_{\pmb{p}}=\hat U_{\pmb{i}_1,\pmb{i}_2}\hat U_{\pmb{i}_2,\pmb{i}_3}\hat U_{\pmb{i}_3,\pmb{i}_4}\hat U_{\pmb{i}_4,\pmb{i}_1} \tag{10.5.6}
\end{equation*}
\begin{figure}[H]
\centering
\includegraphics[width=0.72\textwidth]{fig_10.7.pdf}
\caption*{图 10.7\quad 立方格子中的一条开弦，以及与最近邻链相联系的指标对。开弦的一个端点翻转四个阴影方格上的 $F_{\pmb{p}}$ 的符号。}
\end{figure}
其中 $\sum_{\pmb{p}}$ 对立方格子的所有方格面求和。这里 $\pmb{i}_1$、$\pmb{i}_2$、$\pmb{i}_3$、$\pmb{i}_4$ 标记方格 $\pmb{p}$ 的四个角。引入复费米子算符
\begin{equation*}
2\psi_{1,\pmb{i}}=\lambda_{\pmb{i}}^{x}+\mathrm{i}\lambda_{\pmb{i}}^{\bar{x}}, \qquad 2\psi_{2,\pmb{i}}=\lambda_{\pmb{i}}^{y}+\mathrm{i}\lambda_{\pmb{i}}^{\bar{y}}, \qquad 2\psi_{3,\pmb{i}}=\lambda_{\pmb{i}}^{z}+\mathrm{i}\lambda_{\pmb{i}}^{\bar{z}},
\end{equation*}
并令 $N_{\pmb{i}}=\sum_{a=1,2,3}\psi_{a,\pmb{i}}^{\dagger}\psi_{a,\pmb{i}}$ 为格点 $\pmb{i}$ 上的费米子数算符。证明：在每个格点费米子数目为偶的子空间内，费米子哈密顿量 (10.5.6) 约化为自旋 $3/2$ 哈密顿量 (10.5.4)。
\item 用 10.2 节中的程序，求式 (10.5.4) 在具有周期性边界条件的偶数 $\times$ 偶数 $\times$ 偶数格子上的基态能量与基态简并度。
\end{enumerate}

\subsection{弦算符与闭合弦凝聚}

为了证明 $H_{3D}$ 的基态具有闭合弦凝聚，我们需要找到一个与 $H_{3D}$ 对易的闭合弦算符。首先，让我们构造更普遍的开弦算符。

为了用物理的 $\gamma^{ab}$ 算符构造弦算符，我们注意到，可以给每条最近邻链 $\pmb{i}\to\pmb{j}$ 联系一对指标 $ab$，如图 10.7 所示。例如，链 $\pmb{i}\to\pmb{i}+\pmb{x}$ 对应 $x\bar{x}$，链 $\pmb{i}\to\pmb{i}-\pmb{x}$ 对应 $\bar{x}x$，链 $\pmb{i}\to\pmb{i}+\pmb{y}$ 对应 $y\bar{y}$，等等。开弦
\begin{figure}[H]
\centering
\includegraphics[width=0.72\textwidth]{fig_10.8.pdf}
\caption*{图 10.8\quad 回路 $C$ 可以分成两个回路 $C_1$ 与 $C_2$。}
\end{figure}
算符则定义为
\begin{equation*}
W(C)=-\left(-\mathrm{i}\gamma_{\pmb{i}_1}^{a_1b_1}\right)\left(-\mathrm{i}\gamma_{\pmb{i}_2}^{a_2b_2}\right)\dots\left(-\mathrm{i}\gamma_{\pmb{i}_n}^{a_nb_n}\right) \tag{10.5.7}
\end{equation*}
其中 $b_ma_{m+1}$ 是与链 $\pmb{i}_m\to\pmb{i}_{m+1}$ 相联系的指标对。例如，图 10.7 中那条开弦的开弦算符由下式给出：
\begin{equation*}
W(C)=-\left(-\mathrm{i}\gamma_{\pmb{i}}^{\bar{z}y}\right)\left(-\mathrm{i}\gamma_{\pmb{i}+\pmb{y}}^{\bar{y}x}\right)\left(-\mathrm{i}\gamma_{\pmb{i}+\pmb{x}+\pmb{y}}^{\bar{x}x}\right)
\end{equation*}
我们想强调，开弦 $C$ 是通过相邻格点把诸链的中点连接起来而构成的（见图 10.7）。因此，开弦的端点生活在链上。

闭合弦算符仍由式 (10.5.7) 给出，只是此时 $\pmb{i}_1$ 与 $\pmb{i}_n$ 为最近邻，并且 $b_na_1$ 是与链 $\pmb{i}_n\to\pmb{i}_1$ 相联系的指标对。

让我们先讨论导致闭合弦凝聚的闭合弦算符的一些特殊性质。利用 $\gamma^{ab}$ 的马约拉纳表示，可以证明任意两个闭合弦算符彼此对易。由于 $\hat F_{\pmb{p}}$ 本身就是一个闭合弦算符，闭合弦算符与自旋 $3/2$ 哈密顿量对易。因此，基态也是闭合弦算符的本征态，具有闭合弦凝聚（或更一般地，闭合弦网凝聚）。

如果我们把一个回路 $C$ 分成两个回路 $C_1$ 与 $C_2$（如图 10.8 那样），那么三个回路的闭合弦算符满足关系
\begin{equation*}
W(C)=W(C_1)W(C_2) \tag{10.5.8}
\end{equation*}
这一关系使我们能把 $W(C)$ 写成诸 $F_{\pmb{p}}$ 的乘积，从而计算凝聚闭合弦的振幅。对 $\hat F_{\pmb{p}}=1$ 的基态，我们发现
\begin{equation*}
\langle W(C_{\text{close}})\rangle=1
\end{equation*}
对 $\hat F_{\pmb{p}}=-1$ 的基态，我们发现
\begin{equation*}
\langle W(C_{\text{close}})\rangle=(-)^{N_{\pmb{p}}}
\end{equation*}
其中 $N_{\pmb{p}}$ 是曲面 $S_{C_{\text{close}}}$ 上方格的数目，而 $S_{C_{\text{close}}}$ 是由立方体的面组成、以闭合弦 $C_{\text{close}}$ 为边界的曲面。

\subsection*{问题 10.5.3}

\begin{enumerate}
\item 利用 $\gamma^{ab}=\mathrm{i}\lambda^{a}\lambda^{b}$ 的马约拉纳费米子表示，证明闭合弦算符 (10.5.7) 可以表示为
\begin{equation*}
W(C_{\text{close}})=\hat U_{\pmb{i}_1\pmb{i}_2}\hat U_{\pmb{i}_2\pmb{i}_3}\dots\hat U_{\pmb{i}_{n-1}\pmb{i}_n}\hat U_{\pmb{i}_n\pmb{i}_1}
\end{equation*}
其中 $\hat U_{\pmb{i}\pmb{j}}$ 的定义见式 (10.5.5)。
\item 证明对两个闭合回路定义的闭合弦算符 (10.5.7) 彼此对易。
\item 验证关系式 (10.5.8)。
\end{enumerate}

\subsection{作为开弦端点的演生费米子}

在闭合弦凝聚态中，开弦的端点成为一种新的准粒子。把开弦算符作用到一个基态上，通过计算 $F_{\pmb{p}}$ 与开弦算符之间的对易子可以发现，开弦算符只在其两个端点附近翻转 $F_{\pmb{p}}$ 的符号。因此，弦本身不耗费能量，只有弦端点耗费能量。

注意到弦端点生活在链上，而每条链连接着四个方格，我们发现开弦算符在这四个方格上翻转 $F_{\pmb{p}}$ 的符号（见图 10.7）。因此，开弦的每个端点对应一个小 $Z_2$ 通量圈，耗费能量 $4|2g|=8|g|$。这些小 $Z_2$ 通量圈对应于基态之上的准粒子激发。

为了求出 $Z_2$ 通量圈的统计性质，我们需要用到 Levin 与 Wen (2003) 引入的以下统计跳跃代数（见 4.1.4 节）：
\begin{align*}
&t_{jl}t_{kj}t_{ji}=\mathrm{e}^{\mathrm{i}\theta}t_{ji}t_{kj}t_{jl},\\
&[t_{ij},t_{kl}]=0, \qquad \text{若 } i,j,k,l \text{ 互不相同}, \tag{10.5.9}
\end{align*}
其中 $t_{ji}$ 描述粒子从链 $i$ 到链 $j$ 的跳跃（注意，这里弦的端点生活在链上）。已经证明，若粒子的跳跃算符满足 $\theta=\pi$ 的代数 (10.5.9)，则这些粒子是费米子。如果我们取 $(i,j,k,l)$
\begin{figure}[H]
\centering
\includegraphics[width=0.72\textwidth]{fig_10.9.pdf}
\caption*{图 10.9\quad 链 1 上的一个费米子可以跳跃到十条不同的链 2--11 上。跳跃 $1\to4$ 由 $\gamma_{\pmb{i}}^{zx}$ 产生，跳跃 $1\to3$ 由 $\gamma_{\pmb{i}}^{\bar{x}x}$ 产生，跳跃 $1\to7$ 由 $\gamma_{\pmb{j}}^{y\bar{x}}$ 产生。}
\end{figure}
为图 10.9 中的链 $(2,1,3,4)$，那么跳跃算符 $(t_{jl},t_{kj},t_{ji})$ 由 $(\gamma_{\pmb{i}}^{xz},\gamma_{\pmb{i}}^{\bar{x}x},\gamma_{\pmb{i}}^{x\bar{z}})$ 给出。式 (10.5.9) 成为
\begin{equation*}
\gamma_{\pmb{i}}^{xz}\gamma_{\pmb{i}}^{\bar{x}x}\gamma_{\pmb{i}}^{x\bar{z}}=\mathrm{e}^{\mathrm{i}\theta}\gamma_{\pmb{i}}^{x\bar{z}}\gamma_{\pmb{i}}^{\bar{x}x}\gamma_{\pmb{i}}^{xz}
\end{equation*}
利用 $\gamma^{ab}$ 的马约拉纳费米子表示，我们发现上式在 $\theta=\pi$ 时成立。如果我们取 $(i,j,k,l)$ 为链 $(2,1,3,10)$，则式 (10.5.9) 成为
\begin{equation*}
\gamma_{\pmb{j}}^{\bar{x}x}\gamma_{\pmb{i}}^{\bar{x}x}\gamma_{\pmb{i}}^{x\bar{z}}=\mathrm{e}^{\mathrm{i}\theta}\gamma_{\pmb{i}}^{x\bar{z}}\gamma_{\pmb{i}}^{\bar{x}x}\gamma_{\pmb{j}}^{\bar{x}x}
\end{equation*}
上式同样在 $\theta=\pi$ 时成立。$(i,j,k,l)$ 的其他取法都属于上述两种类型。弦端点的跳跃满足费米子跳跃代数。因此，弦的端点（即小 $Z_2$ 通量圈）是费米子。我们看到，只要玻色模型具有某些类型闭合弦的凝聚，费米子就能从局域玻色模型中演生出来。

\section{量子转子模型与 $U(1)$ 格点规范理论}

\begin{keypoints}
\item 在更高维数，规范理论相当不平凡。它描述一个纠缠态的涨落，而这些涨落无法用局域自由度来表示。
\item $U(1)$ 规范涨落是空间中闭合弦网的涨落。
\item 如果基态包含任意尺寸闭合弦网的凝聚，就可以演生出无能隙的 $U(1)$ 规范玻色子。
\end{keypoints}

我们已经研究过具有演生 $Z_2$ 规范场的精确可解局域玻色模型。本节我们将从量子转子模型（quantum rotor model）出发，研究一个连续 $U(1)$ 规范理论的演生（Foerster et al., 1980; Senthil and Motrunich, 2002; Wen, 2003a）。
\begin{figure}[H]
\centering
\includegraphics[width=0.72\textwidth]{fig_10.10.pdf}
\caption*{图 10.10\quad (a) 四转子系统。(b) 一个简单的格点规范理论，由格点规范场 $a_{i,i+1}$ 与 $a_0(i)$，$i=1,2,3,4$ 描述。}
\end{figure}
我们还将看到，弦网凝聚及相关的量子序如何与演生的 $U(1)$ 规范玻色子相联系。

本节也可以看作是对量子 $U(1)$ 规范理论的一种物理描述。在我看来，用规范势来描述规范理论的标准方式是不自然的，因为规范势本身是非物理的。这里我们将证明，可以用物理自由度——闭合弦网——来表述 $U(1)$ 规范理论（Banks et al., 1977; Savit, 1980）。

\subsection{四转子系统}

\begin{keypoints}
\item 转子角的某些组合的强量子涨落导致一个低能有效规范理论。
\end{keypoints}

让我们先考虑一个单转子：一个质量为 $m$ 的粒子在圆周 $0\leqslant\theta<2\pi$ 上运动。其拉格朗日量为
\begin{equation*}
L=\frac{1}{2}m\dot\theta^{2}+K\cos\theta
\end{equation*}
哈密顿量具有形式
\begin{equation*}
H=\frac{(S^{z})^{2}}{2m}-K\cos\theta=\frac{(S^{z})^{2}}{2m}-\frac{K}{2}(a+a^{\dagger})
\end{equation*}
其中 $S^{z}$ 是角动量，而 $a=\mathrm{e}^{-\mathrm{i}\theta}$。如果选取基底态 $|n\rangle=(2\pi)^{-1/2}\mathrm{e}^{\mathrm{i}n\theta}$（$n$ 取整数），则我们发现 $S^{z}|n\rangle=n|n\rangle$，$a|n\rangle=|n-1\rangle$。

为了得到具有演生低能规范理论的最简单模型，让我们考虑以下由 $\theta_{\langle12\rangle}$、$\theta_{\langle23\rangle}$、$\theta_{\langle34\rangle}$ 和 $\theta_{\langle41\rangle}$ 描述的四转子模型（见图 10.10(a)）：
\begin{equation*}
H=\sum_{i}\left(U\big(S_{\langle i-1,i\rangle}^{z}-S_{\langle i,i+1\rangle}^{z}\big)^{2}+J\big(S_{\langle i,i+1\rangle}^{z}\big)^{2}\right), \tag{10.6.1}
\end{equation*}
其中我们已约定 $4+1\sim1$，$1-1\sim4$。希尔伯特空间由 $|n_{\langle12\rangle}n_{\langle23\rangle}n_{\langle34\rangle}n_{\langle41\rangle}\rangle$ 张成，其中整数 $n_{i,i+1}$ 是 $S_{\langle i,i+1\rangle}^{z}$ 的本征值。若 $U\gg J$，则低能激发由能量 $E=4Jn^{2}$ 的 $|nnnn\rangle$ 态描述。所有其余激发的能量均为 $U$ 的量级。

为了看清它与格点规范理论的联系，我们想写出四转子模型的拉格朗日量。如果把哈密顿量写成 $H=\frac{1}{2}P^{\top}VP$ 的形式，其中 $P^{\top}=(S_{\langle12\rangle}^{z},S_{\langle23\rangle}^{z},S_{\langle34\rangle}^{z},S_{\langle41\rangle}^{z})$，而
\begin{equation*}
V=\begin{pmatrix}
4U+2J & -2U & 0 & -2U\\
-2U & 4U+2J & -2U & 0\\
0 & -2U & 4U+2J & -2U\\
-2U & 0 & -2U & 4U+2J
\end{pmatrix},
\end{equation*}
则拉格朗日量为
\begin{equation*}
L=\frac{1}{2}\dot\Theta^{\top}M\dot\Theta,
\end{equation*}
其中 $\Theta^{\top}=(\theta_{\langle12\rangle},\theta_{\langle23\rangle},\theta_{\langle34\rangle},\theta_{\langle41\rangle})$，$M=V^{-1}$。

显然，在上述拉格朗日量中我们看不到规范理论的任何迹象。为了得到规范理论，我们需要换一种方法导出拉格朗日量。利用 $H$ 的路径积分表示，并注意到 $(\theta,S^{z})$ 是一对正则坐标--动量，我们发现
\begin{equation*}
Z=\int\mathcal{D}S^{z}\,\mathcal{D}\theta\;\mathrm{e}^{\mathrm{i}\int\mathrm{d}t\,\big(\sum_{i}S_{\langle i,i+1\rangle}^{z}\dot\theta_{\langle i,i-1\rangle}-H\big)}
\end{equation*}
引入场 $a_0(i)$（$i=1,2,3,4$）来解耦 $U$ 项，可以把上述路径积分改写为
\begin{equation*}
Z=\int\mathcal{D}S^{z}\,\mathcal{D}\theta\,\mathcal{D}a_0\;\mathrm{e}^{\mathrm{i}\int\mathrm{d}t\,\big(\sum_{i}S_{\langle i,i+1\rangle}^{z}\dot\theta_{\langle i,i+1\rangle}-\tilde H(S^{z},a_0)\big)}
\end{equation*}
其中 $\tilde H=\sum_{i}\Big(J(S_{\langle i,i+1\rangle}^{z})^{2}+a_0(i)\big(S_{\langle i-1,i\rangle}^{z}-S_{\langle i,i+1\rangle}^{z}\big)-\frac{a_0^{2}(i)}{4U}\Big)$。积掉 $S_{\langle i,i+1\rangle}^{z}$ 后，我们得到
\begin{equation*}
Z=\int\mathcal{D}\theta\,\mathcal{D}a_0\;\mathrm{e}^{\mathrm{i}\int\mathrm{d}t\,L(\theta,\dot\theta,a_0)}
\end{equation*}
其中拉格朗日量为
\begin{equation*}
L=\frac{1}{4J}\sum_{i}\left(\big(\dot\theta_{\langle i,i+1\rangle}+a_0(i)-a_0(i+1)\big)^{2}+\frac{a_0^{2}(i)}{4U}\right)
\end{equation*}
在大 $U$ 极限下，可以丢掉 $a_0^{2}(i)/4U$ 项，得到
\begin{equation*}
L=\frac{1}{4J}\sum_{i}\big(\dot a_{i,i+1}+a_0(i)-a_0(i+1)\big)^{2}
\end{equation*}
这正是单个方格上 $U(1)$ 格点规范理论的拉格朗日量，其中
\begin{equation*}
a_{i,i+1}=\theta_{\langle i,i+1\rangle}, \qquad a_{i+1,i}=-\theta_{\langle i,i+1\rangle}
\end{equation*}
作为格点规范场（见图 10.10(b)）。可以检验，上述拉格朗日量在下列变换下不变：
\begin{equation*}
a_{ij}(t)\to a_{ij}(t)+\phi_{j}(t)-\phi_{i}(t), \qquad a_0(i,t)\to a_0(i,t)+\dot\phi_{i}(t) \tag{10.6.2}
\end{equation*}
这个变换称为规范变换。我们注意到，低能波函数 $\Psi(a_{12},a_{23},a_{34},a_{41})$ 是 $|nnnn\rangle$ 态的叠加。所有低能态都是规范不变的，即在规范变换 $a_{ij}\to a_{ij}+\phi_{j}-\phi_{i}$ 下不变。

连续 $U(1)$ 规范理论的电场由 $\pmb{E}=\pmb{\dot a}-\partial a_{0}$ 给出（取 $c=e=1$ 单位）。在格点规范理论中，电场变成定义在链上的如下量：
\begin{equation*}
E_{ij}=\dot a_{ij}-\big(a_0(j)-a_0(i)\big)
\end{equation*}
我们看到，我们的格点规范拉格朗日量可以写成 $L=\frac{1}{4J}\sum_{i}E_{i,i+1}^{2}$。把它与连续 $U(1)$ 规范理论 $\mathcal{L}\propto\pmb{E}^{2}-\pmb{B}^{2}$ 相比较，可见我们的拉格朗日量只包含对应于 $\pmb{E}^{2}$ 的动能。更普遍的格点规范理论还包含对应于 $\pmb{B}^{2}$ 的势能项。

为了得到势能项，我们把转子模型推广为
\begin{equation*}
\begin{split}
H=\sum_{i}\Big(&U\big(S_{\langle i-1,i\rangle}^{z}-S_{\langle i,i+1\rangle}^{z}\big)^{2}+J\big(S_{\langle i,i+1\rangle}^{z}\big)^{2}\\
&+t\big(\mathrm{e}^{\mathrm{i}\theta_{\langle i-1,i\rangle}}\mathrm{e}^{\mathrm{i}\theta_{\langle i,i+1\rangle}}+\text{h.c.}\big)\Big),
\end{split} \tag{10.6.3}
\end{equation*}
我们注意到 $\langle nnn|\mathrm{e}^{\mathrm{i}\theta_{\langle i-1,i\rangle}}\mathrm{e}^{-\mathrm{i}\theta_{\langle i,i+1\rangle}}|nnn\rangle=0$。因此，在 $t$ 的一阶，新项对低能没有影响。新项的低能效应只在 $t$ 的二阶才出现。

如果带着新项重复上述计算，就得到以下拉格朗日量：
\begin{equation*}
L=\frac{1}{4J}\sum_{i}\left(\big(\dot a_{i,i+1}+a_0(i)-a_0(i+1)\big)^{2}-t\big(\mathrm{e}^{\mathrm{i}(a_{i-1,i}+a_{i,i+1})}+\text{h.c.}\big)+\frac{a_0^{2}(i)}{4U}\right)
\end{equation*}
在上述拉格朗日量中，为什么新项在 $t$ 一阶下没有低能效应，稍微难看出来一些。让我们集中考虑如下形式的涨落：\footnote{在格点规范理论中，这种类型的涨落被称为纯规范涨落。}
\begin{equation*}
a_{i-1,i}=\phi_{i}-\phi_{i-1}
\end{equation*}
积掉 $a_0(i)$ 之后，这类涨落的拉格朗日量具有形式 $L=\frac{1}{2}\phi_{i}m_{ij}\phi_{j}-\sum_{i}t\big(\mathrm{e}^{\mathrm{i}(\phi_{i-1}+\phi_{i+1})}+\text{h.c.}\big)$，其中 $m_{ij}=O(U^{-1})$。我们看到，在大 $U$ 极限下，上述形式的涨落是快速而强烈的。由于 $\phi_{i}$ 生活在紧致空间上（即 $\phi_{i}$ 与 $\phi_{i}+2\pi$ 表示同一个点），这些快速涨落全都具有 $U$ 量级的大能隙。现在我们看到，$t$ 项 $t\,\mathrm{e}^{\mathrm{i}(\phi_{i-1}-\phi_{i+1})}$ 由于 $\phi_{i}$ 的强涨落而平均为零，在 $t$ 一阶下没有任何影响。然而，在 $t$ 的二阶，会出现一项 $t^{2}\prod_{i=1,3}\mathrm{e}^{\mathrm{i}(\theta_{\langle i-1,i\rangle}+\theta_{\langle i,i+1\rangle})}=t^{2}\mathrm{e}^{\mathrm{i}\sum_{i}a_{i-1,i}}$。这样的项不依赖于 $\phi_{i}$，因而不会平均为零。于是我们预期低能有效拉格朗日量具有如下形式：
\begin{equation*}
L=\frac{1}{4J}\sum_{i}\left(\big(\dot a_{i,i+1}+a_0(i)-a_0(i+1)\big)^{2}+\frac{a_0^{2}(i)}{4U}\right)+g\cos\Phi \tag{10.6.4}
\end{equation*}
其中 $g=O(t^{2}/U)$，而 $\Phi=\sum_{i}a_{i,i+1}$ 是 $U(1)$ 规范场穿过方格的通量。上述拉格朗日量不依赖快速量子涨落，并且在规范变换 (10.6.2) 下不变。我们看到，正是 $a_{ij}$ 的某些组合的强量子涨落，导致了一个规范不变的低能有效拉格朗日量。

为了定量地计算 $g$，我们先要导出低能有效哈密顿量。如果把 $t$ 项当作微扰来处理，并把低能态当作简并态，那么在 $t$ 的二阶，我们有
\begin{align*}
&\langle n_{1},n_{1},n_{1},n_{1}|H_{eff}|nnnn\rangle\\
&=-\frac{t^{2}}{2U}\langle n_{1},n_{1},n_{1},n_{1}|\mathrm{e}^{\mathrm{i}(\theta_{\langle34\rangle}+\theta_{\langle41\rangle})}|n_{1},n_{1},n,n\rangle\langle n_{1},n_{1},n,n|\mathrm{e}^{\mathrm{i}(\theta_{\langle12\rangle}+\theta_{\langle23\rangle})}|nnnn\rangle\\
&\quad+\text{另外三个类似的项}\\
&=-\frac{2t^{2}}{U}
\end{align*}
其中 $n_{1}=n+1$。于是，低能有效哈密顿量为
\begin{equation*}
H=\sum_{i}\left(U\big(S_{\langle i-1,i\rangle}^{z}-S_{\langle i,i+1\rangle}^{z}\big)^{2}+J\big(S_{\langle i,i+1\rangle}^{z}\big)^{2}\right)-\frac{4t^{2}}{U}\cos(\Phi).
\end{equation*}
相应的拉格朗日量由式 (10.6.4) 给出，其中
\begin{equation*}
g=\frac{4t^{2}}{U}.
\end{equation*}
如前所述，纯规范涨落具有 $U$ 量级的大能隙。$U$ 以下的低能有效理论可以令 $U\to\infty$ 得到，
\begin{figure}[H]
\centering
\includegraphics[width=0.72\textwidth]{fig_10.11.pdf}
\caption*{图 10.11\quad (a) 三转子系统；(b) 相关的格点规范理论。}
\end{figure}
并且我们得到
\begin{equation*}
L=\frac{1}{4J}\sum_{i}\big(\dot a_{i,i+1}+a_0(i)-a_0(i+1)\big)^{2}+g\cos\Phi \tag{10.6.5}
\end{equation*}

这个模型曾在 6.4 节中研究过。那里证明了，当 $g=0$ 时，能级为 $E_{n}=4Jn^{2}$，与式 (10.6.1) 的低能能级精确一致。因此，式 (10.6.1) 在低能下确实是一个规范理论。

这里我们想作一点评论。什么是规范理论？四转子模型在低能下真的是一个规范理论吗？规范理论的标准定义是含有规范势的理论。按照这个定义，如果我们坚持用规范势来写四转子模型的拉格朗日量，它就可以是一个规范理论。另一方面，四转子模型又不是规范理论，因为它的自然拉格朗日量不含有任何规范势。我们看到，规范理论并不是一个定义良好的物理概念。“一个模型是不是规范理论”并不是一个定义良好的问题。这几乎就像：如果我们把四个格点标记为 $(a,b,c,d)$，就管四转子模型叫规范理论；而如果我们把四个格点标记为 $(1,2,3,4)$，就管它叫非规范理论。从这个观点看，本节的讨论是相当形式化的，没有什么物理意义。

\subsection*{问题 10.6.1}

考虑式 (10.6.3) 所描述的 $i=1,2,3$ 三格点转子模型（见图 10.11）。

\begin{enumerate}
\item 导出低能格点规范理论的作用量。你只需把 $g$ 计算到 $O(1)$ 系数的精度。
\item 假定 $U\gg g\gg J$，计算低能下的能级。
\end{enumerate}

\subsection{量子转子格子与人工光}

\begin{keypoints}
\item 量子转子格子可以用一个格点 $U(1)$ 规范理论来描述。
\end{keypoints}

\begin{figure}[H]
\centering
\includegraphics[width=0.72\textwidth]{fig_10.12.pdf}
\caption*{图 10.12\quad 转子格子、表示低能涨落的一个回路，以及一对电荷激发（A、B）。在大 $U$ 极限下，每个阴影方格四角处转子的角动量之和为零。}
\end{figure}

\begin{keypoints}
\item 格点规范理论的解禁闭相包含人工光（artificial light）。
\end{keypoints}

在理解了只含少数几个转子的系统之后，我们就可以着手研究格子转子系统了。作为一个例子，我们将考虑一个正方格子，每条链上有一个转子（见图 10.12）。我们有
\begin{align*}
H=U\sum_{\pmb{i}}\Bigg(\sum_{\pmb{\alpha}}S_{\pmb{i}+\pmb{\alpha}}^{z}\Bigg)^{2}&+J\sum_{\pmb{i},\pmb{\alpha}}\big(S_{\pmb{i}+\pmb{\alpha}}^{z}\big)^{2}+\sum_{\pmb{i}}\Big(t_{1}\,\mathrm{e}^{\mathrm{i}\big(\theta_{\pmb{i}+\frac{\pmb{x}}{2}}-\theta_{\pmb{i}+\frac{\pmb{y}}{2}}\big)}\\
&+t_{2}\,\mathrm{e}^{\mathrm{i}\big(\theta_{\pmb{i}+\frac{\pmb{y}}{2}}-\theta_{\pmb{i}-\frac{\pmb{x}}{2}}\big)}+t_{3}\,\mathrm{e}^{\mathrm{i}\big(\theta_{\pmb{i}-\frac{\pmb{x}}{2}}-\theta_{\pmb{i}-\frac{\pmb{y}}{2}}\big)}+t_{4}\,\mathrm{e}^{\mathrm{i}\big(\theta_{\pmb{i}-\frac{\pmb{y}}{2}}-\theta_{\pmb{i}+\frac{\pmb{x}}{2}}\big)}+\text{h.c.}\Bigg) \tag{10.6.6}
\end{align*}
这里 $\pmb{i}=(i_x,i_y)$ 标记正方格子的格点，而 $\pmb{\alpha}=\pm\frac{\pmb{x}}{2},\pm\frac{\pmb{y}}{2}$。$U$ 项强制施加了“格点 $\pmb{i}$ 附近四个转子的总角动量为零”这一约束。该模型还可以推广到三维立方格子：
\begin{align*}
H=U\sum_{\pmb{i}}\Bigg(&\sum_{\pmb{\alpha}}S_{\pmb{i}+\pmb{\alpha}}^{z}\Bigg)^{2}+J\sum_{\pmb{i},\pmb{\alpha}}\big(S_{\pmb{i}+\pmb{\alpha}}^{z}\big)^{2}\\
&+\sum_{\pmb{i},\langle\pmb{\alpha}_1\pmb{\alpha}_2\rangle}\left(t_{\langle\pmb{\alpha}_1\pmb{\alpha}_2\rangle}\,\mathrm{e}^{\mathrm{i}(\theta_{\pmb{i}+\pmb{\alpha}_1}-\theta_{\pmb{i}+\pmb{\alpha}_2})}+\text{h.c.}\right) \tag{10.6.7}
\end{align*}
这里 $\pmb{i}=(i_x,i_y,i_z)$ 标记立方格子的格点，$\pmb{\alpha}=\pm\frac{\pmb{x}}{2},\pm\frac{\pmb{y}}{2},\pm\frac{\pmb{z}}{2}$，而两个标记 $\pmb{i}+\pmb{\alpha}_1$ 与 $\pmb{i}+\pmb{\alpha}_2$ 表示绕 $\pmb{i}$ 格点的两个最近邻转子。

% 实际覆盖范围（以 PNG 为准）：本块不含 §10.5——§10.5、§10.6.1 及 §10.6.2 前半
% （书页 466–478）属 ch10_c 块；本块自 §10.6.2 中段（书页 479 顶）起，
% 至 §10.7.3 中段（书页 491 末）止。
% 块界说明：
%   起点——书页 479（p494）顶部 "In the following, we are going to show that..."
%   为新起段落，书页 478（p493）末句 "...the two nearest-neighbor rotors around
%   the i site." 已完整结束，故本块不设 %JOIN%，亦无 ch10_c 收尾内容可写入
%   %--C-TAIL-- 区。
%   终点——书页 491（p506）末句 "In the continuum limit, ... and B is the
%   magnetic field" 延续到书页 492（p507，ch10_e 块）顶部 "strength of the
%   U(1) gauge field. Comparing this with..."。按约定，半句译文（"……而 $\pmb B$ 是"）
%   留在本文件正文最末，未写入任何 TAIL 区；请 ch10_e 以 %JOIN% 起头续译
%   "该 $U(1)$ 规范场的磁场强度。把它与……"，勿重复勿遗漏。

下面我们将证明，在 $U\gg t$、$J$ 且 $t^{2}/U\gg J$ 的极限下，上述二维和三维模型都包含一个低能集体模式。由于其非局域的性质，这种集体模式与自旋波、声子这类通常的集体模式非常不同。实际上，它不对应于任何局域序参量的涨落，必须用 $U(1)$ 规范理论来描写。我们把这样的集体模式称为 $U(1)$ 规范涨落。该集体模式在二维模型中具有指数式小的能隙 $\Delta\sim\mathrm{e}^{-C\sqrt{t^{2}/JU}}$，在三维模型中则严格无能隙。三维模型中的集体模式在各个方面都表现得像我们宇宙中的光。因此，我们也可以称它为人工光。

上述模型引人注目之处在于，它们在格子尺度上没有任何规范结构的迹象。$U(1)$ 规范涨落及与之相联系的规范结构是在低能下演生的。这提示了一种可能：我们宇宙中的光或许也以类似的方式出现。规范结构可能会在强关联量子系统中自然地演生。本节要研究的转子模型是包含人工光的最简单模型。下面为简单起见，我们将集中于二维转子模型。计算与结果可以容易地推广到三维转子模型。

为了理解我们二维转子系统的动力学，让我们先设 $J=t=0$ 且 $U>0$。此时，哈密顿量由执行局域投影的对易项构成。基态高度简并，构成一个低能子空间。其中一个基态是每个转子的 $S^{z}_{\pmb i}=0$ 的态。其他基态可以从这个基态出发来构造：在正方格子上画一条有向回路，然后沿回路交替地把转子的角动量增加或减少 1（见图 10.12）。重复这一过程可以构造出所有简并基态。我们看到，投影空间中的涨落由回路表示。尽管低能子空间是通过局域投影得到的，它却带有一些非局域的特征。若 $t$ 与 $J$ 非零，则 $t$ 项会使这些回路涨落，而 $J$ 项会给这些回路正比于回路长度的能量。显然，当 $U\gg J,t$ 时，系统的低能性质由回路的涨落决定。系统可以有两个不同的相。当 $J\gg t$ 时，系统只有少许小回路涨落。当 $J\ll t$ 时，系统有许多大回路涨落，它们甚至可以充满整个空间。后面我们会看到，这些大回路涨落实际上正是 $U(1)$ 规范涨落。我们注意到，回路可以相交与交叠。实际上，一个典型的回路看起来更像一张闭合弦的网。下面，为强调这种涨落的分支结构，我们把回路涨落称为闭合弦网涨落。

简并基态在由
\begin{equation*}
U(\phi_{\pmb i})=\mathrm{e}^{\mathrm{i}\sum_{\pmb i}((-)^{\pmb i}\phi_{\pmb i}\sum_{\alpha}S^{z}_{\pmb i+\alpha})}\tag{10.6.8}
\end{equation*}
生成的局域对称变换下不变，其中 $(-)^{\pmb i}=(-)^{i_x+i_y}$。上述变换正是规范变换。因此，我们也可以说简并基态是规范不变的。

用与 10.6.1 小节相似的计算，我们发现我们的格子转子模型在大 $U$ 极限下可以用如下低能有效拉格朗日量描写：
\begin{equation*}
L=\frac{1}{4J}\sum_{\pmb i,\pmb\mu=\pmb x,\pmb y}[\dot{a}_{\pmb i,\pmb i+\pmb\mu}+a_0(\pmb i)-a_0(\pmb i+\pmb\mu)]^{2}+g\sum_{\pmb i}\cos(\Phi_{\pmb i})\tag{10.6.9}
\end{equation*}
这里 $a_{\pmb i,\pmb i+\pmb x}=(-)^{\pmb i}\theta_{\pmb i+\frac{\pmb x}{2}}$，$a_{\pmb i,\pmb i+\pmb y}=(-)^{\pmb i}\theta_{\pmb i+\frac{\pmb y}{2}}$，而 $a_{\pmb i\pmb j}=-a_{\pmb j\pmb i}$。另外，$\Phi_{\pmb i}$ 是通过 $\pmb i$ 处方格的 $U(1)$ 通量，即 $\Phi_{\pmb i}=a_{\pmb i,\pmb i+\pmb x}+a_{\pmb i+\pmb x,\pmb i+\pmb x+\pmb y}+a_{\pmb i+\pmb x+\pmb y,\pmb i+\pmb y}+a_{\pmb i+\pmb y,\pmb i}$。推广 10.6.1 小节的二阶微扰计算，我们得到
\begin{equation*}
g=2(t_1t_3+t_2t_4)/U.
\end{equation*}

式 (10.6.9) 描写一个 $U(1)$ 格点规范理论。

这里我们想强调，具有 $U(1)$ 规范理论形式的拉格朗日量 (10.6.9) 并不一定意味着低能激发的行为像规范玻色子。这是因为量子涨落可以非常强，来自拉格朗日量的直观图像可能完全造成误导。要有行为像规范玻色子的激发，我们需要选取 $J$ 与 $g$，使拉格朗日量 (10.6.9) 处于半经典极限。为了看清何时能够达到半经典极限，我们把格点规范理论的作用量写成如下无量纲形式：
\begin{equation*}
S=\sqrt{\frac{g}{J}}\int\mathrm{d}\tilde t\left(\frac{1}{4}\sum_{\pmb i,\pmb\mu=\pmb x,\pmb y}[\partial_{\tilde t}a_{\pmb i,\pmb i+\pmb\mu}+\tilde a_0(\pmb i)-\tilde a_0(\pmb i+\pmb\mu)]^{2}+\sum_{\pmb i}\cos(\Phi_{\pmb i})\right)
\end{equation*}
其中 $\tilde t=\sqrt{gJ}\,t$，$\tilde a_0=a_0/\sqrt{gJ}$。我们发现，当 $g/J=t^{2}/JU\gg1$ 时即达到半经典极限。在这一极限中，我们的模型包含 $U(1)$ 规范玻色子（即人工光）作为其唯一的低能集体激发。

由于瞬子效应，$U(1)$ 规范激发在 $1+2$ 维中会产生能隙（Polyakov, 1977）（见 6.3.2 小节）。瞬子效应与一个方格上 $U(1)$ 通量 $\Phi$ 从 0 到 $2\pi$ 的改变相联系。为了估计瞬子效应的重要性，让我们考虑只含单个方格的模型（即前面讨论过的四转子模型）。这样的模型由式 (6.4.11) 描写。瞬子效应对应于一条路径 $\Phi(t)$（依赖时间的通量），其中 $\Phi$ 从 $\Phi(-\infty)=0$ 变到 $\Phi(+\infty)=2\pi$。为了估计瞬子作用量，我们设
\begin{equation*}
\Phi(t)=\left\{\begin{array}{ll}0,&\text{当 } t<0\\ 2\pi t/T,&\text{当 } 0<t<T\\ 2\pi,&\text{当 } T<t\end{array}\right.
\end{equation*}
瞬子作用量的最小值为
\begin{equation*}
S_c=\pi\sqrt{g/J}
\end{equation*}
此时 $T=\pi/2\sqrt{gJ}$。由瞬子气体在时间轴上的密度 $\sqrt{Jg}\,\mathrm{e}^{-S_c}$，我们估计 $U(1)$ 规范玻色子的能隙为
\begin{equation*}
\Delta_{\mathrm{gauge}}\sim\sqrt{Jg}\,\mathrm{e}^{-\pi\sqrt{g/J}}\tag{10.6.10}
\end{equation*}

我们看到，当 $g/J\gg1$ 时，能隙可以非常小，低能涨落非常像无能隙的光子。当然，真正无能隙的人工光只存在于 $(3+1)$ 维模型中，例如三维转子模型 (10.6.7)。

三维转子模型在大 $U$ 极限下具有如下低能有效拉格朗日量：
\begin{equation*}
L=\frac{1}{4J}\sum_{\pmb i,\pmb\mu=\pmb x,\pmb y,z}[\dot{a}_{\pmb i,\pmb i+\pmb\mu}+a_0(\pmb i)-a_0(\pmb i+\pmb\mu)]^{2}+\sum_{\pmb p}g_{\pmb p}\cos(\Phi_{\pmb p})\tag{10.6.11}
\end{equation*}
这里 $\sum_{\pmb p}$ 对所有方格求和，$\Phi_{\pmb p}$ 是通过方格 $\pmb p$ 的 $U(1)$ 通量。另外，$g_{\pmb p}$ 可以依赖方格的取向，其数值可通过式 (10.6.7) 中的 $t_{\langle\pmb\alpha_1\pmb\alpha_2\rangle}$ 调节。

\subsection{人工光与人工电荷的弦网理论}

\begin{keypoints}
\item 规范理论是披着伪装的闭合弦网理论。
\item 禁闭相是具有稀疏小弦的相；库仑相是具有充满整个空间的大闭合弦网的相。
\item 规范玻色子（在库仑相中）对应于大闭合弦网的涨落，而规范电荷是开弦的端点。
\end{keypoints}

如前所述，我们的转子模型中 $U$ 以下的低能激发由 $S^{z}$ 增/减的闭合弦网描写。为了把这个图像弄得更精确，我们想在格子上定义一个弦网理论。

弦网理论的希尔伯特空间是我们转子模型 (10.6.6) 的希尔伯特空间的一个子空间。为了构造弦网希尔伯特空间，我们首先需要引入一个弦产生算符。弦产生算符由 $\mathrm{e}^{\pm\mathrm{i}\theta_{\langle\pmb i\pmb j\rangle}}$ 算符的乘积构成：
\begin{figure}[H]
\centering
\includegraphics[width=0.72\textwidth]{fig_10.13.pdf}
\caption*{图 10.13\quad (a) 闭合弦网；(b) 开弦网。}
\end{figure}
\begin{equation*}
U(C)=\prod_{\langle\pmb i\pmb j\rangle}\mathrm{e}^{\mathrm{i}(-)^{\pmb i}\theta_{\langle\pmb i\pmb j\rangle}}\tag{10.6.12}
\end{equation*}
其中 $C$ 是连接正方格子上最近邻格点的一根弦，乘积 $\prod_{\langle\pmb i\pmb j\rangle}$ 遍历构成该弦的所有最近邻链 $\langle\pmb i\pmb j\rangle$，而 $\theta_{\langle\pmb i\pmb j\rangle}$ 是链 $\langle\pmb i\pmb j\rangle$ 上转子的 $\theta$ 场。

由于弦可以相交与交叠，弦看起来更像一张网。我们把 $C$ 称为弦网，把 $U(C)$ 称为弦网算符。若 $C$ 由一组闭合弦构成，我们就称它为闭合弦网；若 $C$ 至少含有一根开弦，就称它为开弦网（见图 10.13）。

弦网希尔伯特空间包含一个所有 $S^{z}_{\pmb I}=0$ 的态。按定义，这样的态对应于无弦态。如果我们把闭合弦算符 (10.6.12) 作用到 $S^{z}_{\pmb I}=0$ 态上，就得到弦网希尔伯特空间中的另一个态。这样的态由沿闭合弦 $C$ 的 $S^{z}_{\pmb I}=\pm1$ 构成，可以视为含一根闭合弦 $C$ 的态。弦网希尔伯特空间中的其他态对应于多弦态，可以通过反复把闭合弦算符 (10.6.12) 作用到 $S^{z}_{\pmb I}=0$ 态上产生。

我们弦网理论的哈密顿量由限制在弦网子空间中的转子哈密顿量 (10.6.6) 得到，其形式为
\begin{equation*}
H_{str}=\sum_{\langle\pmb i\pmb j\rangle}J(S^{z}_{\pmb i\pmb j})^{2}-\sum_{\pmb p}\frac{1}{2}(gW_{\pmb p}+h.c.)\tag{10.6.13}
\end{equation*}
其中求和 $\sum_{\langle\pmb i\pmb j\rangle}$ 遍历正方格子上所有最近邻链 $\langle\pmb i\pmb j\rangle$，$S^{z}_{\pmb i\pmb j}$ 是最近邻链 $\langle\pmb i\pmb j\rangle$ 上转子的角动量算符，$\sum_{\pmb p}$ 遍历正方格子的所有方格，而 $W_{\pmb p}$ 是绕方格 $\pmb p$ 的闭合弦的闭合弦算符。可以验证，上述哈密顿量确实作用在弦网希尔伯特空间之内。$J$ 项赋予弦网中的弦有限的弦张力，而 $g$ 项作为弦的“跳跃”项使弦网发生涨落。式 (10.6.13) 中的 $J$ 项直接来自式 (10.6.6) 中的 $J$ 项，而 $g$ 项来自式 (10.6.6) 中的 $t$ 项的二阶微扰展开。

由构造可以清楚地看到，弦网希尔伯特空间与我们的模型 (10.6.6) 的低能希尔伯特空间（由能量小于 $U$ 的态构成）恒同。由有效格点规范理论 (10.6.9) 的推导也同样清楚，弦网哈密顿量 (10.6.13) 与格点规范拉格朗日量 (10.6.9) 直接相关。实际上，格点规范理论的哈密顿量与弦网哈密顿量 (10.6.13) 恒同。弦网理论中的 $\sum_{\pmb i\pmb j}J(S^{z}_{\pmb i\pmb j})^{2}$ 项对应于规范理论中的 $\frac{1}{4J}\sum_{\langle\pmb i\pmb j\rangle}[\dot{a}_{\pmb i\pmb j}+a_0(\pmb i)-a_0(\pmb j)]^{2}$ 项，弦网理论中的 $\sum_{\pmb p}\frac{1}{2}g(W_{\pmb p}+h.c.)$ 项对应于规范理论中的 $g\sum_{\pmb p}\cos(\Phi_{\pmb p})$ 项。由于 $S^{z}\sim J^{-1}\dot{\theta}_{\langle\pmb i\pmb j\rangle}=J^{-1}\dot{a}_{\pmb i\pmb j}$ 对应于沿链的电通量，$S^{z}$ 增/减的闭合弦便对应于电通量管的一个回路。我们想强调，规范理论与弦网理论之间的上述联系不仅适用于二维模型 (10.6.6)，也适用于任何维数中的模型。

我们看到，$U(1)$ 规范理论 (10.6.9) 实际上是闭合弦网的一个动力学理论。通常人们期望闭合弦网的动力学理论像 (10.6.13) 那样用弦网写成。然而，由于我们更熟悉场论，前几节所做的事情可以视为用场论描写弦网理论的一种尝试。借助某些数学技巧，我们达到了目的：能够把弦网理论写成规范场论的形式。规范场论是一类特殊的场论，其中的场\emph{并不}对应于物理自由度，而且其物理希尔伯特空间是非局域的（即总物理希尔伯特空间不能写成局域希尔伯特空间的直积）。我们在这里想说明的要点是：规范理论（至少这里讨论的规范理论）是一种披着伪装的闭合弦网理论。换言之，规范理论与闭合弦网理论互为对偶。实际上，如果把时间离散化、考虑一个时空格子，我们就能找到 $U(1)$ 格点规范理论与时空中的膜的统计模型之间的精确映射（Banks et al., 1977; Savit, 1980）。

在大 $J/g$ 极限（因而是大 $\Delta_{\mathrm{gauge}}$ 极限）下，转子模型与弦网模型的基态都是每个转子 $S^{z}=0$ 的态。在这个相中，闭合弦网（或电通量管）涨落不大，能量正比于其长度。这意味着 $U(1)$ 规范理论处于禁闭相。在小 $J/g$ 极限下，闭合弦网强烈涨落，空间充满任意尺寸的闭合弦网。根据上一小节的计算，我们注意到小 $J/g$ 相也可以视为具有无能隙规范玻色子的库仑相。把这两幅图像结合起来，我们看到无能隙规范玻色子对应于大闭合弦网的涨落。

在把闭合弦网与人工光联系起来之后，我们现在转向人工电荷（artificial charge）。要为人造 $U(1)$ 规范场产生一对带相反人工电荷的粒子，我们需要画一根开弦，并沿弦交替地增减转子的 $S^{z}$（见图 10.12）。开弦的端点，作为电通量管的端点，对应于带相反人工电荷的粒子。我们注意到，与转子不同，带电粒子生活在正方格子的格点上。在禁闭相中，连接两个人工电荷的弦涨落不大，弦的能量正比于弦的长度。我们看到，人工电荷之间存在线性禁闭。

在小 $J/g$ 极限下，大的 $g$ 引起闭合弦网的强烈涨落，从而导致无能隙的 $U(1)$ 规范涨落。连接两个电荷的弦的强烈涨落还把线性禁闭势变为电荷之间的 $\log(r)$ 势。

为了理解带人工电荷的粒子的动力学并导出 $\log(r)$ 势，让我们导出这些带电粒子的低能有效理论。让我们先设 $J=t=0$。把开弦算符 (10.6.12) 作用到基态上，可以产生一对带相反单位人工电荷的带电粒子。在开弦的端点处，$(\sum_{\alpha}S^{z}_{\pmb i+\alpha})^{2}=1$。我们发现每个带电粒子具有能量 $U$，它来自 $U\sum_{\pmb i}(\sum_{\alpha}S^{z}_{\pmb i+\alpha})^{2}$ 这一项。弦本身不耗费能量。如果把同一开弦算符作用 $n$ 次，就在开弦的端点产生 $n$ 个单位的相反人工电荷。$n$ 个单位电荷的能量是 $n^{2}U$。

含电荷的规范理论包含两部分，即电荷与弦。让我们先只考虑电荷，把带电粒子当作独立粒子处理。此时，带电粒子的总希尔伯特空间由态 $|\{n_{\pmb i}\}\rangle$ 构成，其中 $n_{\pmb i}$ 是正方格子格点 $\pmb i$ 上人工电荷的数目。这里 $|\{n_{\pmb i}\}\rangle$ 是一个能量本征态，能量为 $E=U\sum_{\pmb i}n^{2}_{\pmb i}$。这样的系统可以用拉格朗日量
\begin{equation*}
L=\sum_{\pmb i}\frac{1}{4U}\dot{\varphi}^{2}_{\pmb i}\tag{10.6.14}
\end{equation*}
描写，其中 $\varphi_{\pmb i}$ 是一个角变量。产生单位电荷带电粒子的算符是 $\mathrm{e}^{\mathrm{i}\varphi_{\pmb i}}$。

规范涨落由弦和式 (10.6.9) 描写。在 $U(1)$ 规范理论中，基态是规范不变的（在远离电荷处）。由式 (10.6.8) 可见，规范不变性意味着 $\sum_{\alpha}S^{z}_{\pmb i+\alpha}=0$。因此在远离电荷处没有弦的开端。

为了找到电荷与规范场的耦合理论，我们把带电粒子总是处在开弦端点这一事实包括进来。换言之，规范理论的物理态包含开弦，但开弦只能终止在电荷上。这一约束可以由规范不变性施加。于是我们可以利用规范不变性，把电荷拉格朗日量 (10.6.14) 与规范拉格朗日量 (10.6.9) 合在一起，写出耦合理论。利用规范不变性，我们发现组合拉格朗日量具有如下形式（见问题 6.4.5）：
\begin{equation*}
L=\sum_{\pmb i}\frac{(\dot{\varphi}_{\pmb i}+a_0(\pmb i))^{2}}{4U}+\sum_{\pmb i,\pmb\mu=\pmb x,\pmb y}\frac{[\dot{a}_{\pmb i,\pmb i+\pmb\mu}+a_0(\pmb i)-a_0(\pmb i+\pmb\mu)]^{2}}{4J}+g\sum_{\pmb i}\cos(\Phi_{\pmb i})
\end{equation*}
把规范场包括进来之后，单电荷产生算符 $\mathrm{e}^{\mathrm{i}\varphi_{\pmb i}}$ 不再是物理的，因为它不是规范不变的。规范不变的算符
\begin{equation*}
\mathrm{e}^{-\mathrm{i}\varphi_{\pmb i_1}}\mathrm{e}^{\mathrm{i}a_{\pmb i_1\pmb i_2}}\mathrm{e}^{\mathrm{i}a_{\pmb i_2\pmb i_3}}\dots\mathrm{e}^{\mathrm{i}a_{\pmb i_{N-1}\pmb i_N}}\mathrm{e}^{\mathrm{i}\varphi_{\pmb i_N}}
\end{equation*}
总是产生一对相反电荷，连同连接两电荷的开弦。因此，开弦总是终止在电荷上。实际上，上述规范不变算符正是开弦算符 (10.6.12)。我们还看到，弦算符 (10.6.12) 与 Wegner--Wilson 圈算符密切相关（Wegner, 1971; Wilson, 1974; Kogut, 1979）。

$t$ 项产生带电粒子在正方格子上次近邻之间的跳跃。因此，若 $t\neq0$，带电粒子将具有非平庸的色散。相应的拉格朗日量为
\begin{equation*}
L=\sum_{\pmb i}\frac{(\dot{\varphi}_{\pmb i}+a_0(\pmb i))^{2}}{4U}-\sum_{\langle\pmb i\pmb j\rangle,a=1,2}t(\mathrm{e}^{\mathrm{i}(\varphi_{\pmb i}-\varphi_{\pmb j}-a_{\pmb i\pmb k_a}-a_{\pmb k_a\pmb j})}+h.c.)\tag{10.6.15}
\end{equation*}
其中 $\langle\pmb i\pmb j\rangle$ 是正方格子上的次近邻，$\pmb k_{1,2}$ 是格点 $\pmb i$ 与 $\pmb j$ 之间的两个格点。上述拉格朗日量还告诉我们，带电粒子是玻色子。

我们还注意到，$S^{z}_{\pmb i\pmb j}$ 的增加对应于 $\pmb i$ 与 $\pmb j$ 处的两个人工电荷。因此，每一单位人工电荷携带 1/2 角动量！（注意，总角动量 $\sum_{\langle\pmb i\pmb j\rangle}S^{z}_{\langle\pmb i\pmb j\rangle}$ 是守恒量。）

\subsection{二维与三维转子系统的物理性质}

\begin{keypoints}
\item 格点规范理论的连续极限。
\item 人工光的速度与精细结构常数。
\end{keypoints}

为了理解二维模型中人工光的物理性质，让我们按如下方式取式 (10.6.9) 的连续极限：
\begin{equation*}
a_{\pmb i\pmb j}=\delta\pmb x_{\langle\pmb i\pmb j\rangle}\cdot\pmb a(\pmb x),\qquad a_0(\pmb i)=a_0(\pmb x),
\end{equation*}
其中 $\pmb a=(a_x,a_y)$ 是一个二维矢量场（二维中的矢量规范势），$a_0$ 对应于势场，$\pmb x$ 取在格点 $\pmb i$ 附近，$\delta\pmb x_{\langle\pmb i\pmb j\rangle}$ 是连接正方格子上 $\pmb i$ 与 $\pmb j$ 两格点的矢量，而 $l$ 是正方格子的晶格常数。在连续极限下，拉格朗日量 (10.6.9) 变为
\begin{equation*}
L=\int\mathrm{d}^{2}\pmb x\ \left(\frac{1}{4J}\pmb e^{2}-\frac{gl^{2}}{2}\pmb b^{2}\right)\tag{10.6.16}
\end{equation*}
其中 $\pmb e=\partial_t\pmb a-\partial_{\pmb x}a_0$ 与 $\pmb b=\partial_x a_y-\partial_y a_x$ 分别是相应的仿造电场与仿造磁场。我们看到，人工光的速度为 $c_a=\sqrt{2gJl^{2}/h^{2}}$，人工光的带宽约为 $E_a=2c_a\hbar/l=\sqrt{8gJ}$。

由式 (10.6.15)，我们得到描写二维模型中带电粒子（在 $U\gg t$ 极限下）的连续拉格朗日量：
\begin{align*}
L=\int\mathrm{d}^{2}\pmb x\ \sum_{I=1,2}\Big(\phi^{\dagger}_{I}&(\mathrm{i}\partial_t-a_0-U)\phi_I-\frac{8tl^{2}}{2}\big|(\partial_{i}+\mathrm{i}a_{i})\phi_I\big|^{2}\\
+\bar{\phi}^{\dagger}_{I}&(\mathrm{i}\partial_t+a_0-U)\bar{\phi}_I-\frac{8tl^{2}}{2}\big|(\partial_{i}-\mathrm{i}a_{i})\bar{\phi}_I\big|^{2}\Big)
\end{align*}
（译注：原文两处时间导数下标印作 $\partial_l$，应为 $\partial_t$，此处已按本意译出。）

其中 $\phi_I$ 描写带正电的玻色子，$\bar{\phi}_I$ 描写带负电的玻色子，$\psi_1$ 与 $\bar{\psi}_1$ 描写正方格子偶格点上带电的玻色子，$\psi_2$ 与 $\bar{\psi}_2$ 描写奇格点上带电的玻色子。产生一对带电玻色子耗费能量 $2U$。玻色子的质量为 $m=(8tl^{2})^{-1}$，且 $mc_a^{2}=gJ/4t$。我们想指出，玻色子速度 $\sim 8tl$ 可以大于人工光的速度。正电荷与负电荷之间的势能为 $V(r)=\frac{2J}{\pi}\ln r$。

对于三维模型，连续极限下的拉格朗日量为
\begin{equation*}
L=\int\mathrm{d}^{3}\pmb x\ \left(\frac{1}{4Jl}\pmb e^{2}-\frac{gl}{2}\pmb b^{2}\right)\tag{10.6.17}
\end{equation*}
其中 $\pmb e$ 与 $\pmb b$ 分别是三维中的仿造电场与仿造磁场。人工光的速度为 $c_a=\sqrt{2gJl^{2}/h^{2}}$。上述三维拉格朗日量可以改写成更标准的形式
\begin{equation*}
L=\int\mathrm{d}^{3}\pmb x\ \frac{1}{8\pi\alpha}\left(\frac{1}{c_a}\pmb e^{2}-c_a\pmb b^{2}\right)\tag{10.6.18}
\end{equation*}
其中 $\alpha=\frac{1}{2\pi}\sqrt{J/2g}$ 是人工精细结构常数。带电玻色子的质量为 $m=(8tl^{2})^{-1}$。人工原子（artificial atom，即一个带正电玻色子与一个带负电玻色子的束缚态）具有 $\frac{1}{4}mc_a^{2}\alpha^{2}$ 量级的能级间距和 $2/\alpha mc_a$ 量级的尺寸。

\subsection*{问题 10.6.2}

验证式 (10.6.16) 与 (10.6.17)。

\section{从 $SU(N_f)$ 自旋模型演生的光与电子}

\begin{keypoints}
\item 局域玻色模型中非平庸的量子序（即弦网凝聚）可以同时导致无质量的光子与无质量的费米子。弦网凝聚提供了统一光与费米子的一条途径。
\item 光子与费米子的无质量性由刻画弦网凝聚的 PSG 保护。
\end{keypoints}

正如本章开头所强调的，理解光与电子的存在，就是理解这些粒子为什么是无质量的（与 Planck 质量相比近乎无质量）。对于一个一般的相互作用系统，无质量（即无能隙）激发是罕见的。如果它们存在，那么它们的存在必有原因。

我们知道，对称性破缺能够产生并保护无能隙的 Nambu--Goldstone 模。9.10 节中提出，除对称性破缺之外，量子序及与之相联系的 PSG 也能产生并保护无能隙激发。由量子序产生并保护的无能隙激发可以是无能隙规范玻色子和/或无能隙费米子。

在本节中，我们将研究立方格子上一个具体的 $SU(N_f)$ 自旋模型（Wen, 2002a）。该模型的基态具有非平庸的量子序（即弦网凝聚）。结果，模型具有演生的无质量 $U(1)$ 规范玻色子（人工光）以及无质量的带电费米子（人工电子与质子）。换言之，$SU(N_f)$ 自旋模型具有一个演生的 QED！基于这种 $SU(N_f)$ 自旋模型的世界将与我们自己的世界十分相似；两个世界都由通过库仑相互作用耦合的带电费米子（即电子与质子）构成。更复杂的立方格子自旋模型甚至可以导致含有光子、胶子、轻子和夸克的演生 QED 与 QCD（Wen, 2003b）。因此，如果有一天我们发现自然界中所有基本粒子都从一个推广的自旋模型（即局域玻色模型）中演生出来，那也不足为奇。

\subsection{立方格子上的 $SU(N_f)$ 自旋模型}

该模型由 $SU(N_f)$ 自旋构成，立方格子的每个格点上放有一个自旋。这里 $N_f$ 为偶数，格点用三个整数标记，即 $\pmb i=(i_x,i_y,i_z)$。每个格点上的态
\begin{equation*}
|a_1,a_2,\dots,a_{N_f/2}\rangle,\qquad a_1,\dots,a_{N_f/2}=1,2,\dots,N_f
\end{equation*}
构成 $SU(N_f)$ 的秩为 $N_f/2$ 的反对称张量表示。每个格点上的自旋算符 $S^{a\bar b}_{\pmb i}$（$a,\bar b=1,2,\dots,N_f$）构成 $SU(N_f)$ 的伴随表示：
\begin{gather*}
[S^{a\bar b}_{\pmb i},S^{c\bar d}_{\pmb i}]=\delta^{\bar b c}S^{a\bar d}_{\pmb i}-\delta^{a\bar d}S^{c\bar b}_{\pmb i}\\
S^{a\bar b}_{\pmb i}\delta_{a\bar b}=0\\
(S^{a\bar b}_{\pmb i})^{\dagger}=S^{\bar b a}_{\pmb i}\tag{10.7.1}
\end{gather*}
其中重复指标意味着求和。

令
\begin{equation*}
W_{\pmb i}=-N_f^{-4}S^{ab}_{\pmb i+\pmb x}S^{bc}_{\pmb i+\pmb x+\pmb y}S^{cd}_{\pmb i+\pmb y}S^{da}_{\pmb i}
\end{equation*}
我们把 $SU(N_f)$ 模型的哈密顿量取为（Wen, 2002a）
\begin{equation*}
H=g\sum_{\pmb i}[W_{\pmb i}+h.c]\tag{10.7.2}
\end{equation*}

\subsection{$SU(N_f)$ 模型的基态}

为理解基态，让我们引入 $SU(N_f)$ 自旋模型的费米子表示（Affleck and Marston, 1988）。我们先引入构成 $SU(N_f)$ 基础表示的 $N_f$ 个费米子算符 $\psi_a$。每个格点上的希尔伯特空间于是由具有 $N_f/2$ 个费米子的态构成。自旋算符由下式给出：
\begin{equation*}
S^{a\bar b}=\psi_a\psi^{\dagger}_{\bar b}-\frac{1}{2}\delta_{a\bar b}
\end{equation*}
如果我们引入
\begin{equation*}
\hat\chi_{\pmb i\pmb j}=N_f^{-1}\psi^{\dagger}_{a\pmb i}\psi_{a\pmb j}
\end{equation*}
则 $W_{\pmb i}$ 由下式给出：
\begin{equation*}
W_{\pmb i}=\hat\chi_{\pmb i,\pmb i+\pmb x}\hat\chi_{\pmb i+\pmb x,\pmb i+\pmb x+\pmb y}\hat\chi_{\pmb i+\pmb x+\pmb y,\pmb i+\pmb y}\hat\chi_{\pmb i+\pmb y,\pmb i}+O(1/N_f)
\end{equation*}
利用费米子表示，我们可以把哈密顿量 (10.7.2) 写成
\begin{equation*}
H=g\sum_{\pmb p}\Big[\prod_{\square}\hat\chi_{\pmb i\pmb j}+h.c.\Big]+O(N_f^{-1})\tag{10.7.3}
\end{equation*}
其中 $\sum_{\pmb p}$ 对所有方格 $\pmb p$ 求和，$\prod_{\square}$ 对方格 $\pmb p$ 周围的四条链求积；也就是说，$\prod_{\square}\hat\chi_{\pmb i\pmb j}\allowbreak=\hat\chi_{\pmb i_1\pmb i_2}\allowbreak\hat\chi_{\pmb i_2\pmb i_3}\allowbreak\hat\chi_{\pmb i_3\pmb i_4}\allowbreak\hat\chi_{\pmb i_4\pmb i_1}$，其中 $\pmb i_1,\pmb i_2,\pmb i_3,\pmb i_4$ 是方格 $\pmb p$ 周围的四个格点。

我们注意到，在大 $N_f$ 极限下，$\hat\chi_{\pmb i\pmb j}$ 彼此对易，即 $[\hat\chi_{\pmb i\pmb j},\hat\chi_{\pmb i'\pmb j'}]=O(1/N_f)$。另外，$\hat\chi_{\pmb i\pmb j}$ 是 $N_f$ 项之和。尽管每一项都有很强的量子涨落，其和却几乎没有涨落。因此，在大 $N_f$ 极限下 $\hat\chi_{\pmb i\pmb j}$ 的行为像一个 c 数。也就是说，对一个给定的态，我们可以把 $\hat\chi_{\pmb i\pmb j}$ 换成一个复数 $\chi_{\pmb i\pmb j}$。这样的态的能量为 $E=g\sum_{\pmb p}[\prod_{\square}\chi_{\pmb i\pmb j}+h.c.]+O(N_f^{-1})$。不同的态有不同的 $\chi_{\pmb i\pmb j}$，因而有不同的能量。

$SU(N_f)$ 模型的基态可以这样得到：选取一组使 $E=g\sum_{\pmb p}[\prod_{\square}\chi_{\pmb i\pmb j}+h.c.]$ 极小的 $\chi_{\pmb i\pmb j}$。这里我们遇到一个困难：$\chi_{\pmb i\pmb j}$ 并不独立，我们不能任意选取它们的值。例如，我们不能让所有 $\chi_{\pmb i\pmb j}$ 同时为零。我们将用 9.1.1 小节中的投影构造\footnote{9.1.1 小节中的讨论针对的是 $N_f=2$ 的 $SU(N_f)$ 模型。}来克服这一困难。我们从平均场哈密顿量
\begin{equation*}
H_{\mathrm{mean}}=-\sum_{\langle\pmb i\pmb j\rangle}\tilde\chi^{\dagger}_{\pmb i\pmb j}\psi^{\dagger}_{\pmb i}\psi_{\pmb j}+\sum_{\pmb i}a_0(\pmb i)\psi^{\dagger}_{\pmb i}\psi_{\pmb i}
\end{equation*}
出发。我们这样选取 $a_0(\pmb i)$，使得 $H_{\mathrm{mean}}$ 的平均场基态，即 $|\Psi^{\{\tilde\chi_{\pmb i\pmb j}\}}_{\mathrm{mean}}\rangle$，平均每格点拥有 $N_f/2$ 个费米子。$SU(N_f)$ 模型的态可以由平均场态 $|\Psi^{\{\tilde\chi_{\pmb i\pmb j}\}}_{\mathrm{mean}}\rangle$ 投影到每格点恰有 $N_f/2$ 个费米子的子空间中来构造：
\begin{equation*}
|\Psi^{\{\tilde\chi_{\pmb i\pmb j}\}}\rangle=\mathcal P|\Psi^{\{\tilde\chi_{\pmb i\pmb j}\}}_{\mathrm{mean}}\rangle
\end{equation*}
现在我们可以选取 $\tilde\chi_{\pmb i\pmb j}$，使平均能量 $\langle\Psi^{\{\tilde\chi_{\pmb i\pmb j}\}}|H|\Psi^{\{\tilde\chi_{\pmb i\pmb j}\}}\rangle$ 极小，从而得到基态的尝试能量与尝试波函数。

在大 $N_f$ 极限下，投影 $\mathcal P$ 只使波函数 $|\Psi^{\{\tilde\chi_{\pmb i\pmb j}\}}_{\mathrm{mean}}\rangle$ 发生很小的改变。因此，为得到基态能量，我们可以转而做一件更容易的计算：求 $\langle\Psi^{\{\tilde\chi_{\pmb i\pmb j}\}}_{\mathrm{mean}}|H|\Psi^{\{\tilde\chi_{\pmb i\pmb j}\}}_{\mathrm{mean}}\rangle$ 的极小。

极小化按如下步骤进行。我们先挑出一组 $\tilde\chi_{\pmb i\pmb j}$，然后选取 $a_0(\pmb i)$ 使得 $\langle\Psi^{\{\tilde\chi_{\pmb i\pmb j}\}}_{\mathrm{mean}}|\psi^{\dagger}_{\pmb i}\psi_{\pmb i}|\Psi^{\{\tilde\chi_{\pmb i\pmb j}\}}_{\mathrm{mean}}\rangle=N_f/2$，接着计算 $\chi_{\pmb i\pmb j}=N_f^{-1}\langle\Psi^{\{\tilde\chi_{\pmb i\pmb j}\}}_{\mathrm{mean}}|\psi^{\dagger}_{\pmb i}\psi_{\pmb j}|\Psi^{\{\tilde\chi_{\pmb i\pmb j}\}}_{\mathrm{mean}}\rangle$。这些 $\chi_{\pmb i\pmb j}$ 就是上面讨论过的算符 $\hat\chi_{\pmb i\pmb j}$ 所对应的 c 数。我们看到，$\tilde\chi_{\pmb i\pmb j}$ 是独立的，其值可以任意选取。非独立的 $\chi_{\pmb i\pmb j}$ 是 $\tilde\chi_{\pmb i\pmb j}$ 的函数，即 $\chi_{\pmb i\pmb j}(\{\tilde\chi_{\pmb i\pmb j}\})$。基态通过在改变 $\tilde\chi_{\pmb i\pmb j}$ 的同时使 $E=g\sum_{\pmb p}[\prod_{\square}\chi_{\pmb i\pmb j}(\{\tilde\chi_{\pmb i\pmb j}\})+h.c.]$ 极小而得到。

我们发现，基态可以通过如下选取 $\tilde\chi_{\pmb i\pmb j}$ 得到：
\begin{align*}
\tilde\chi_{\pmb i,\pmb i+\hat{\pmb x}}&\equiv-\mathrm{i}, & \tilde\chi_{\pmb i,\pmb i+\hat{\pmb y}}&\equiv-\mathrm{i}(-)^{i_x},\\
\tilde\chi_{\pmb i,\pmb i+\hat{\pmb z}}&\equiv-\mathrm{i}(-)^{i_x+i_y}, & a_0(\pmb i)&=0
\end{align*}
它引导出如下一组 $\chi_{\pmb i\pmb j}$：
\begin{align*}
\bar\chi_{\pmb i,\pmb i+\hat{\pmb x}}&\equiv-\mathrm{i}|\chi|, & \bar\chi_{\pmb i,\pmb i+\hat{\pmb y}}&\equiv-\mathrm{i}(-)^{i_x}|\chi|,\\
\bar\chi_{\pmb i,\pmb i+\hat{\pmb z}}&\equiv-\mathrm{i}(-)^{i_x+i_y}|\chi| & a_0(\pmb i)&=0.\tag{10.7.4}
\end{align*}
拟设 $\bar\chi_{\pmb i\pmb j}$ 在每个方格中都有 $\pi$ 通量。结果是 $\prod_{\square}\chi_{\pmb i\pmb j}=-|\chi|^{4}$。正是这个负号使上述 $\chi_{\pmb i\pmb j}$ 极小化 $g\sum_{\pmb p}[\prod_{\square}\chi_{\pmb i\pmb j}+h.c.]$（假定 $g>0$）。

由拟设 $\bar\chi_{\pmb i\pmb j}$ 描写的基态不破坏任何对称性。然而，该基态具有非平庸的量子序。下面我们将证明，基态支持无质量的 $U(1)$ 规范玻色子与无质量的带电费米子。

\subsection{$SU(N_f)$ 模型的低能动力学}

$\chi_{\pmb i\pmb j}$ 的涨落描写我们 $SU(N_f)$ 自旋模型中的一个集体模式。这一集体模式的动力学由立方格子上的一个经典场论描写。$|\chi|$ 的涨落对应于一种有质量的激发，可以忽略。集体模式的低能涨落由围绕 $\bar\chi_{\pmb i\pmb j}$ 的相位涨落描写，即 $\chi_{\pmb i\pmb j}=\bar\chi_{\pmb i\pmb j}\mathrm{e}^{\mathrm{i}a_{\pmb i\pmb j}}$。如 9.1.1 小节所讨论的，由 $a_{\pmb i\pmb j}$ 描写的涨落是 $U(1)$ 规范涨落。$U(1)$ 规范涨落的物理波函数由变形拟设的投影平均场态给出：
\begin{equation*}
\mathcal P|\Psi^{\{\bar\chi_{\pmb i\pmb j}\mathrm{e}^{\mathrm{i}a_{\pmb i\pmb j}}\}}_{\mathrm{mean}}\rangle.\tag{10.7.5}
\end{equation*}

把上式代入式 (10.7.3)，就得到 $a_{\pmb i\pmb j}$ 的如下有效哈密顿量：
\begin{equation*}
H=-2g\sum_{\pmb p}|\chi|^{4}\cos(\Phi_{\pmb p})\tag{10.7.6}
\end{equation*}
其中
\begin{equation*}
\Phi_{\pmb p}=a_{\pmb i_1\pmb i_2}+a_{\pmb i_2\pmb i_3}+a_{\pmb i_3\pmb i_4}+a_{\pmb i_4\pmb i_1}
\end{equation*}
而 $(\pmb i_1,\pmb i_2,\pmb i_3,\pmb i_4)$ 是方格 $\pmb p$ 周围的四个格点。哈密顿量 (10.7.6) 描写一个 $U(1)$ 格点规范理论：$a_{\pmb i\pmb j}$ 是格点 $U(1)$ 规范场，$\Phi_{\pmb p}$ 是通过方格 $\pmb p$ 的 $U(1)$ 通量。

$U(1)$ 规范涨落 (10.7.5) 不是唯一类型的低能激发。第二类低能激发由带有若干粒子--空穴激发的投影平均场态给出，例如
\begin{equation*}
\mathcal P\psi^{a}_{\pmb i_1}\psi^{b\dagger}_{\pmb i_2}|\Psi^{\{\bar\chi_{\pmb i\pmb j}\}}_{\mathrm{mean}}\rangle.\tag{10.7.7}
\end{equation*}
这些激发对应于带电费米子。

为确定费米子激发的动力学，我们写出
\begin{equation*}
\hat\chi_{\pmb i\pmb j}=\bar\chi_{\pmb i\pmb j}+N_f^{-1}:\psi^{\dagger}_{a,\pmb i}\psi_{a,\pmb j}:\tag{10.7.8}
\end{equation*}
其中 $\bar\chi_{\pmb i\pmb j}=\langle\hat\chi_{\pmb i\pmb j}\rangle$。把式 (10.7.8) 代入式 (10.7.3)，我们得到
\begin{equation*}
H_f=gN_f^{-1}\sum_{\pmb i\pmb j}[\chi'_{\pmb i\pmb j}:\psi^{\dagger}_{\pmb j}\psi_{\pmb i}:+h.c.]+\sum_{\pmb i}a_0(\pmb i):\psi^{\dagger}_{\pmb i}\psi_{\pmb i}:+O\big((:\psi^{\dagger}\psi:)^{2}\big)\tag{10.7.9}
\end{equation*}
其中
\begin{gather*}
\chi'_{\pmb i,\pmb i+\hat{\pmb x}}\equiv4\mathrm{i}|\chi|^{3},\\
\chi'_{\pmb i,\pmb i+\hat{\pmb y}}\equiv4\mathrm{i}(-)^{i_x}|\chi|^{3},\\
\chi'_{\pmb i,\pmb i+\hat{\pmb z}}\equiv4\mathrm{i}(-)^{i_x+i_y}|\chi|^{3}.\tag{10.7.10}
\end{gather*}
我们看到，式 (10.7.9) 描写 $\pi$ 通量相中费米子的跳跃。把式 (10.7.6) 与 (10.7.9) 相结合，就得到费米子与规范激发的下述低能有效哈密顿量：
\begin{align*}
H_{eff}=gN_f^{-1}\sum_{\pmb i\pmb j}[\chi'_{\pmb i\pmb j}\mathrm{e}^{\mathrm{i}a_{\pmb i\pmb j}}:\psi^{\dagger}_{\pmb j}\psi_{\pmb i}:+h.c.]&+\sum_{\pmb i}a_0(\pmb i):\psi^{\dagger}_{\pmb i}\psi_{\pmb i}:\\
&-2g\sum_{\pmb p}|\chi|^{4}\cos(\Phi_{\pmb p})\tag{10.7.11}
\end{align*}
在连续极限下，$-2g\sum_{\pmb p}|\chi|^{4}\cos(\Phi_{\pmb p})$ 变为 $\int\mathrm{d}^{3}\pmb x\,(3l_0 g)|\chi|^{4}\pmb B^{2}$，其中 $l_0$ 是立方格子的晶格常数，而 $\pmb B$ 是
% ----（书页 491 末句延续到书页 492 顶部 "strength of the U(1) gauge field. Comparing this with..."。
% ---- 半句译文止于"而 $\pmb B$ 是"，请 ch10_e 以 %JOIN% 起头续译"该 $U(1)$ 规范场的磁场强度。
% ---- 把它与……"，勿重复勿遗漏。
该 $U(1)$ 规范场的磁场强度。把它与 $U(1)$ 规范场论的标准哈密顿量，即 $H=g_1\pmb B^{2}+g_2\pmb E^{2}$（其中 $\pmb E$ 是 $U(1)$ 规范场的电场）相比较，我们发现我们的 $U(1)$ 哈密顿量缺少来自电场的动能项。动能项以我们模型中 $1/N_f$ 展开的高阶项出现。

为估计 $\pmb E^{2}$ 项中系数的值，我们注意到，对与时间无关的规范势，$\pmb E^{2}$ 具有 $(\partial_{\pmb x}a_0)^{2}$ 的形式。这样的项可以由式 (10.7.11) 积掉费米子而产生。产生的项量级为 $N_f\frac{1}{gN_f^{-1}|\chi|}(l_0\partial_{\pmb x}a_0)^{2}$。第一个系数 $N_f$ 来自如下事实：$N_f$ 族费米子中的每一族对 $(l_0\partial_{\pmb x}a_0)^{2}$ 的贡献都相同。第二个系数描写来自其中一族费米子的贡献。该系数具有能量倒数的量纲，其值可以用费米子带宽的倒数来估计，即 $1/(gN_f^{-1}|\chi|)$。于是，$U(1)$ 规范场的拉格朗日量具有如下形式（Marston and Affleck, 1989; Wen, 2002a）：
\begin{equation*}
\mathcal L=\frac{CN_f^{2}}{gl_0}\pmb E^{2}-(3l_0g|\chi|^{4})\pmb B^{2}
\end{equation*}
其中 $C$ 是一个 $O(1)$ 常数。这样的拉格朗日量描写我们 $SU(N_f)$ 自旋模型中的一种人工光。把上式与标准拉格朗日量 (10.6.18) 相比较，我们发现人工光的速度为 $c_a\sim l_0g/N_f$ 量级，而精细结构常数为 $1/N_f$ 量级。

在动量空间中，费米子跳跃哈密顿量 (10.7.9) 具有如下形式（注意对基态有 $a_0=0$）：
\begin{equation*}
H_f=\sum_{\pmb k}^{\prime}\Psi^{\dagger}_{a,\pmb k}\Gamma(\pmb k)\Psi_{a,\pmb k}
\end{equation*}
其中
\begin{gather*}
\Psi^{\top}_{a,\pmb k}=(\psi_{a,\pmb k},\psi_{a,\pmb k+\pmb Q_x},\psi_{a,\pmb k+\pmb Q_y},\psi_{a,\pmb k+\pmb Q_x+\pmb Q_y}),\\
\pmb Q_x=(\pi,0,0),\qquad \pmb Q_y=(0,\pi,0),\\
\Gamma(\pmb k)=-8|\chi|N_f^{-1}\big(\sin(k_x)\Gamma_1+\sin(k_y)\Gamma_2+\sin(k_z)\Gamma_3\big)\tag{10.7.12}
\end{gather*}
且 $\Gamma_1=\tau^{3}\otimes\tau^{0}$，$\Gamma_2=\tau^{1}\otimes\tau^{3}$，$\Gamma_3=\tau^{1}\otimes\tau^{1}$。这里 $\tau^{1,2,3}$ 是泡利矩阵，$\tau^{0}$ 是 $2\times2$ 单位矩阵。动量求和 $\sum_{\pmb k}^{\prime}$ 的取值范围为 $k_x\in(-\pi/2,\pi/2)$、$k_y\in(-\pi/2,\pi/2)$、$k_z\in(-\pi,\pi)$。由于 $\{\Gamma_i,\Gamma_j\}=2\delta_{ij}$（$i,j=1,2,3$），我们发现费米子具有色散
\begin{equation*}
E(\pmb k)=\pm8g|\chi|^{3}N_f^{-1}\sqrt{\sin^{2}(k_x)+\sin^{2}(k_y)+\sin^{2}(k_z)}
\end{equation*}
我们看到，该色散在 $\pmb k=0$ 与 $\pmb k=(0,0,\pi)$ 处有两个节点。因此，在连续极限下，式 (10.7.9) 将给出 $2N_f$ 个无质量的四分量狄拉克费米子。

在纳入 $U(1)$ 规范涨落之后，无质量的狄拉克费米子作为带单位电荷的费米子与 $U(1)$ 规范场相互作用。因此，我们 $SU(N_f)$ 自旋模型的总有效理论是一个具有 $2N_f$ 族单位电荷狄拉克费米子的 QED。我们将把这些费米子称为人工电子。连续有效理论具有如下形式：
\begin{equation*}
\mathcal L=\dot\psi_{I,a}D_0\gamma^{0}\psi_{I,a}+v_f\dot\psi_{I,a}D_i\gamma^{i}\psi_{I,a}+\frac{CN_f^{2}}{gl_0}\pmb E^{2}-(3l_0g|\chi|^{4})\pmb B^{2}+\dots\tag{10.7.13}
\end{equation*}
其中 $I=1,2$，$D_0=\partial_t+\mathrm{i}a_0$，$D_i=\partial_i+\mathrm{i}a_i|_{i=1,2,3}$，$v_f=8l_0g|\chi|^{3}/N_f^{-1}$，$\gamma_\mu|_{\mu=0,1,2,3}$ 是 $4\times4$ 狄拉克矩阵，且 $\bar\psi_{I,a}=\psi^{\dagger}_{I,a}\gamma^{0}$。这里 $\psi_{1,a}$ 与 $\psi_{2,a}$ 是构成 $SU(N_f)$ 基础表示的狄拉克费米子场。我们想指出，虽然人工光的速度 $c_a$ 与人工电子的速度 $v_f$ 都是 $l_0g/N_f$ 量级，但在我们的模型中这两个速度不必相同。因此，洛伦兹对称性并无保证。（译注：式 (10.7.13) 与 $v_f$ 均按原书排印转写；按此处 $\bar\psi_{I,a}$ 的定义，式中前两个 $\psi$ 上的点号疑为 $\bar\psi$ 横杠之误印，$v_f$ 表达式中的除号按色散关系疑为乘号之误，即 $v_f=8l_0g|\chi|^{3}N_f^{-1}$。）

式 (10.7.13) 描写 $SU(N_f)$ 模型在一个量子有序相——$\pi$ 通量相——中的低能动力学。费米子与规范玻色子都是无质量的，并且彼此相互作用。这里我们想讨论一个重要问题：在积掉高能的费米子涨落与规范涨落之后，费米子与规范玻色子还保持无质量吗？一般而言，无质量激发之间的相互作用会产生质量项，除非这种无质量性受到对称性或其他某种机制的保护。对我们的 $SU(N_f)$ 模型来说，基态不破坏任何对称性，所以我们不能用自发破缺的对称性来解释这些无质量激发。无质量激发受到刻画基态中量子序（或弦网凝聚）的 PSG 的保护（见 9.10 节及 Wen (2002a)）。

\subsection{评注：关于规范理论与费米统计的一些历史评注}

\begin{keypoints}
\item 看待规范场有两种方式：把它视为局域相位不变性的几何对象，或视为一个关联系统的集体模式。
\item “规范”的含义。
\item 规范场和费米场的存在并不意味着规范玻色子和费米子会作为低能准粒子出现。
\end{keypoints}

第一个系统的规范理论是 Maxwell 的电磁理论。虽然人们为了表示电场和磁场而引进了矢量势 $A_\mu$，但 $A_\mu$ 的含义并不清楚。

规范场的概念由 Weyl 于 1918 年引进，他还提出矢量势 $A_\mu$ 是一个规范场。Weyl 的想法受到 Einstein 引力理论的启发，是统一电磁学与引力的一次尝试。在 Einstein 的广义相对论中，坐标不变性导致了引力。于是 Weyl 想到，另一种几何对象的不变性或许会导致电磁学。他提出的是标度不变性。

考虑一个取值为 $f$ 的物理量。我们知道，若不指明单位，数值 $f$ 本身是没有意义的。用 $\omega$ 表示单位，这个物理量实际上由 $f\omega$ 给出。这就是标度的相对性。现在假设该物理量在空间每一点都有定义（这样我们考虑的就是一个物理场）。我们想知道，如何比较不同点 $x^{\mu}$ 与 $x^{\mu}+\mathrm{d}x^{\mu}$ 处的物理量。我们不能只比较数值 $f(x^{\mu})$ 与 $f(x^{\mu}+\mathrm{d}x^{\mu})$，因为单位 $\omega$ 在不同的点可能不同。对邻近的点 $x^{\mu}$ 与 $x^{\mu}+\mathrm{d}x^{\mu}$，两个单位仅相差一个接近 1 的因子，我们可以把这样的因子写成 $1+S_\mu\mathrm{d}x^{\mu}$。在 $x^{\mu}$ 与 $x^{\mu}+\mathrm{d}x^{\mu}$ 两处物理量之差不是由 $f(x^{\mu}+\mathrm{d}x^{\mu})-f(x^{\mu})=\partial_\mu f\,\mathrm{d}x^{\mu}$ 给出，而是由 $f(x^{\mu}+\mathrm{d}x^{\mu})(1+S_\mu\mathrm{d}x^{\mu})-f(x^{\mu})=(\partial_\mu+S_\mu)f\,\mathrm{d}x^{\mu}$ 给出。Weyl 证明，局域标度不变性要求只有 $S_\mu$ 的旋度才具有物理意义，正如在 Maxwell 理论中只有 $A_\mu$ 的旋度才有意义一样。这样，Weyl 把 $S_\mu$ 认作矢量势 $A_\mu$。Weyl 把局域标度不变性称为“Eich Invarianz”，这个词被译成了“规范不变性”（gauge invariance）。

然而，Weyl 的想法是错误的，矢量势 $A_\mu$ 并不能被认作那个“规范场” $S_\mu$。不过，Weyl 已经几乎对了。如果我们把物理场看作一个复波函数\footnote{复波函数的概念于 1925 年引进，比 Weyl 的“规范理论”晚了七年。}的振幅，把单位 $\omega$ 看作一个复相位，即 $|\omega|=1$，那么不同点处振幅之差便由 $(\partial_\mu+\mathrm{i}S_\mu)f\,\mathrm{d}x^{\mu}$ 给出，其中不同点的单位相差一个因子 $(1+\mathrm{i}S_\mu\mathrm{d}x^{\mu})$。正是这样的 $S_\mu$ 才能被认作矢量势。所以 $A_\mu$ 实际上应该叫作“相位场”，而“规范不变性”实际上应该叫作“相位不变性”。然而，旧名称一直沿用至今。

这段历史是试图赋予非物理的矢量 $A_\mu$ 某种物理（或几何）意义的一次尝试。它把矢量势视为纤维丛的联络。这一图像已被广泛接受。我们如今把矢量势称为规范场，把 Maxwell 理论称为规范理论。然而，这并不意味着我们必须从局域相位不变性出发把矢量势解释成一个几何对象。毕竟，量子波函数的相位是非物理的。

关于规范理论还有另一种观点。理论物理学界的许多思考者对规范势 $A_\mu$ 的冗余性颇为不满。早在 20 世纪 70 年代初人们就认识到，可以用规范不变的圈算符来刻画规范理论的不同相（Wegner, 1971; Wilson, 1974; Kogut and Susskind, 1975）。后来人们发现，可以用闭合弦来表述整个规范理论（Banks et al., 1977; Foerster, 1979; Gliozzi et al., 1979; Mandelstam, 1979; Polyakov, 1979; Savit, 1980）。这些研究揭示了规范理论与闭合弦理论之间的密切关系——这一观点与矢量势的几何观念大不相同。

在凝聚态物理的一项相关进展中，人们发现，如果玻色模型处于某些量子相，规范场可以从局域玻色模型中演生出来。这一现象也称为规范场的动力学产生。规范场从局域玻色模型中的演生有一段漫长而复杂的历史。演生的 $U(1)$ 规范场最早是在 $(1+1)$ 维 $CP^{N}$ 模型的量子无序相中引进的（D'Adda et al., 1978; Witten, 1979）。在凝聚态物理中，人们在正方格子上玻色自旋模型的自旋液体态的从属玻色子方法中发现了 $U(1)$ 规范场（Affleck and Marston, 1988; Baskaran and Anderson, 1988）。从属玻色子方法不仅含有 $U(1)$ 规范场，还含有无能隙的费米子场。然而，由于瞬子效应及其导致的 $U(1)$ 规范场在 $1+1$ 维和 $1+2$ 维中的禁闭（Polyakov, 1975），上述规范场与无能隙费米子场没有一个能给出作为低能物理准粒子出现的规范玻色子和无能隙费米子。即使在可以忽略瞬子效应的大 $N$ 极限下，$2+1$ 维中 $U(1)$ 规范场与无质量狄拉克费米子之间的边缘耦合也会破坏费米子传播子和规范传播子中的准粒子极点。这曾导致一种看法：$U(1)$ 规范场与无能隙费米子场只不过是“不可靠”的从属玻色子方法的一种非物理的人为产物。因此，寻找演生规范玻色子与演生费米子的关键，不在于写下一个含有规范场和费米场的拉格朗日量，而在于证明规范玻色子和费米子确实出现在物理的低能谱中。实际上，对任何给定的物理系统，我们总可以随意设计一个带有任意所选规范场的拉格朗日量来描写它。然而，拉格朗日量中的规范场未必给出一个作为低能准粒子出现的规范玻色子。只有当规范场的动力学使规范场处于解禁闭相时，规范玻色子才能作为低能准粒子出现。因此，在 D'Adda et al. (1978)、Witten (1979)、Baskaran and Anderson (1988) 与 Affleck and Marston (1988) 的最初发现之后，许多研究者一直在努力寻找规范场的解禁闭相。

在高能物理中，Foerster et al. (1980) 构造了一个具有演生的解禁闭 $U(1)$ 规范玻色子的 $(3+1)$ 维局域玻色模型。该工作提出，自然界中的光可能是演生的。在凝聚态物理中，有人证明，如果在一个二维自旋 1/2 模型中破坏时间反演对称性，那么由于 Chern--Simons 项的出现，从属玻色子方法中的 $U(1)$ 规范场可以处于解禁闭相（Khveshchenko and Wiegmann, 1989; Wen et al., 1989）。这一解禁闭相对应于自旋 1/2 模型的一个自旋液体态（Kalmeyer and Laughlin, 1987），称为手征自旋液体。第二种解禁闭相是通过把 $U(1)$ 规范结构破缺到 $Z_2$ 规范结构而找到的。这种相含有一个解禁闭的 $Z_2$ 规范理论（Read and Sachdev, 1991; Wen, 1991a），称为 $Z_2$ 自旋液体（或短程 RVB 态）。手征自旋液体和 $Z_2$ 自旋液体都有一些惊人的性质。准粒子激发携带自旋 1/2，对应于一次自旋翻转的\emph{一半}。尽管我们的自旋 1/2 模型是纯玻色模型，这些准粒子还可以携带分数统计或费米统计。这些凝聚态的例子表明，规范场和费米统计都可以从局域玻色模型中演生出来。

我们想指出，自旋液体并不是局域玻色模型演生费米子的第一个例子。演生费米子——或更一般地，演生任意子——的第一个例子由 FQH 态给出。虽然 Arovas et al. (1984) 只讨论了任意子如何从磁场中的费米子系统演生出来，但同样的论证很容易推广，用以表明费米子和任意子如何从磁场中的玻色子系统演生出来。此外，1987 年，在共振价键（RVB）态的研究中，人们在正方格子上的最近邻二聚体模型中提出了演生费米子（自旋子）（Kivelson et al., 1987; Rokhsar and Kivelson, 1988; Read and Chakraborty, 1989）。然而，按照解禁闭的图像，Kivelson et al. (1987) 与 Rokhsar and Kivelson (1988) 的结果只在二聚体模型的基态处于 $Z_2$ 解禁闭相时才成立。看来，只含最近邻二聚体的正方格子上的二聚体液体并不是解禁闭态（Rokhsar and Kivelson, 1988; Read and Chakraborty, 1989），因此，正方格子上的最近邻二聚体模型（Rokhsar and Kivelson, 1988）是否具有费米型准粒子仍不清楚（Read and Chakraborty, 1989）。不过，在三角格子上，二聚体液体确实是一个 $Z_2$ 解禁闭态（Moessner and Sondhi, 2001）。因此，Kivelson et al. (1987) 与 Rokhsar and Kivelson (1988) 的结果对三角格子二聚体模型成立，费米型准粒子确实在三角格子上的二聚体液体中演生出来。

上面所有具有演生费米子的模型都是 $(2+1)$ 维模型，在其中，演生费米子可以通过把通量束缚到带电粒子上来理解（Arovas et al., 1984）。最近，Levin and Wen (2003) 指出，演生费米子的关键是一种弦结构。在任意维数中，费米子一般都可以作为开弦的端点出现。这一弦图像使我们能够构造一个具有演生费米子的 $(3+1)$ 维局域玻色模型（Levin and Wen, 2003）。既然规范玻色子与费米子都可以作为弦网凝聚的结果而演生，我们可以说，弦网凝聚提供了统一规范玻色子与费米子的一条途径。

把二维正方格子上的玻色 $SU(N)$ 自旋模型（Affleck and Marston, 1988）加以推广，人们发现，从一个三维立方格子上的玻色 $SU(N)$ 自旋模型中，无能隙的解禁闭 $U(1)$ 规范玻色子和无能隙费米子都演生出来了（Wen, 2002a）。在 $1+3$ 维中，这两种无能隙激发可以彼此分离，因为它们在低能下相互作用很弱。$U(1)$ 规范玻色子和无能隙费米子在各个方面都表现得像光子和电子。因此，这个玻色 $SU(N)$ 自旋模型不仅包含人工光，还包含人工电子。

在规范理论和费米统计诞生约一百年之后，我们正面临如下问题。规范场的起源是什么——是几何的，还是动力学的？费米统计的起源是什么——是与生俱来的，还是演生的？在本书中，我们支持规范玻色子与费米子的动力学演生起源。规范玻色子和费米统计也许只不过是量子多玻色子系统的集体现象，仅此而已。
% ---- 第10章及全书正文至此结束（书页 496 末"...and nothing more."）。
% ---- 其后为 BIBLIOGRAPHY（书页 497–501）与 INDEX（书页 502–504），不属正文翻译范围。
